% 本文件由 LaTeX 排版 Agent 根据有效 translation.json 自动生成。
% author=Minucius Felix; work_id=0410; title=奥克塔维乌斯

\chapter{\WorkTitle{奥克塔维乌斯}}

% structure: p0001 source=h1 target=omitted reason=page-title-already-rendered-by-parent

% structure: p0002 source=h2 target=omitted reason=duplicate-page-title

当我思考并在心中回顾我对\Person{奥克塔维乌斯}的记忆，我那卓越而最忠实的同伴，他的甜美和魅力如此紧附于我，以至于我在某种程度上似乎回到了过去的时光，而不仅仅是回忆早已过去的事情。因此，实际注视他从我眼前消失的程度，却深植于我的心中和最私密的感情中。那位卓越而神圣的人，当他离开\Emph{此生}时，给我留下了对他无限的怀念，这并非无因，尤其是因为他自己也始终对我怀有如此的爱，以致无论是在娱乐还是事务中，他都在意愿上与我一致，同样喜欢或同样厌恶相同的事物。你会以为我们两人共享一个心灵。因此，他是我爱情中的唯一知己，是我错误中的同伴；而当阴霾散去，我从黑暗的深渊中浮现到智慧与真理的光明中时，他并没有抛弃他的同伴，反而——更荣耀的是——超越了他。因此，当我的思绪穿越我们亲密和友谊的全部时期时，我心灵的指向主要固定在他的那次论述上，在那次论述中，他以极有力的论据将仍然执著于迷信虚妄的\Person{凯基利乌斯}，皈依了真正的宗教。\par

% structure: p0004 source=h2 target=omitted reason=duplicate-page-title

因为，为着事务和探望我，\Person{奥克塔维乌斯}匆忙赶到\Place{罗马}，离开了他的家、他的妻子、他的孩子，以及孩子中最吸引人的部分——当他们天真的岁月还在尝试仅说出半句话时，那种因口齿不清的不完美而更加甜美的语言。对于他的到来，我无法用言语表达我多么巨大而急切的喜悦，因为一个如此亲爱的人的意外出现大大增加了我的快乐。因此，一两天后，当我们频繁的相聚满足了情感的渴望，并通过相互叙述得知了因分离而不了解的彼此一切后，我们同意前往那座非常宜人的城市\Place{奥斯蒂亚}，以便我的身体能从海水浴中获得舒缓和适合的干燥体液的治疗，尤其是因为葡萄收获季节的法院假期使我从忧虑中解脱出来。因为那时，在夏日之后，秋季正趋向更温和的温度。因此，当我们在清晨沿着\Place{台伯河}岸走向大海时，既为了呼吸的空气轻轻清爽我们的四肢，也为了松软的沙地在轻松的脚步下极度愉快地下沉；\Person{凯基利乌斯}看到一尊\OriginalTerm{塞拉皮斯}{Serapis}像，按照迷信平民的习俗，将手举到嘴边，用嘴唇吻了它一下。\par

% structure: p0006 source=h2 target=omitted reason=duplicate-page-title

于是\Person{奥克塔维乌斯}说：\Quote{善人的本分，我的兄弟\Person{马尔库斯}，并非如此抛弃一个无论在家或在外都守在你身边的人，在如此粗俗的无知黑暗中，以致你容忍他在这样光天化日之下投身于石头——尽管它们被雕刻成图像、被涂抹膏油、被戴上冠冕——因为你深知，他这错误的耻辱，对你的信誉的损害并不亚于对他自己的损害。} 带着他这番话，我们穿过了城市与海之间的路程，现在我们正行走在宽阔开阔的海岸上。那里，轻轻涟漪的波浪正在抚平外面的沙滩，仿佛要为散步平整地面；由于大海总是动荡不安，即使风平浪静时，它也涌上沙滩，虽然没有泡沫和浪花，却有皱起的卷波。就在那时，我们非常喜欢它的变幻，在水边，我们正弄湿脚底，波浪时而靠近拍打我们的脚，时而退去，溯回原路，把自己吸回自身。就这样，我们缓慢而安静地走着，沿着微微弯曲的海岸，用故事消磨时间。这些故事由\Person{奥克塔维乌斯}讲述，他正在谈论航海。但当我们用谈话占用了散步相当合理的时间后，我们再次折返原路，沿着原路返回。当我们来到那个地方，那里的小船被拖上橡木架，静止地停放着，免于朽烂之险，我们看到一些男孩急切地做手势，玩着往海里扔贝壳的游戏。这游戏是：从海岸上选一个贝壳，被波浪抛掷得磨得光滑；用手指水平捏住贝壳；沿着倾斜且尽可能低的角度在波浪上旋转它，使得扔出时要么掠过波浪背面，要么平滑地滑行游动，要么切入浪顶跳跃起来，仿佛被反复弹簧抬起。那个男孩声称胜利，他的贝壳飞得最远，并且跳跃次数最多。\par

% structure: p0008 source=h2 target=omitted reason=duplicate-page-title

这样，当我们大家都沉浸在这场竞赛的欢乐中时，\Person{Cæcilius}却毫不在意，也不对比赛发笑；他沉默不安，独自站在一旁，由他的表情承认他正因我不知道的什么事而悲伤。我遂问他：\Quote{怎么回事？\Person{Cæcilius}，为什么我认不出你平时的活泼？为什么我徒然寻找那即使在严肃场合也属于你眼神的欢快？}于是他回答说：\Quote{一段时间以来，我们的朋友\Person{Octavius}的讲话使我深感痛苦和不安，他攻击你，责备你疏忽，以便借这个指控更严厉地谴责我的无知。因此我将更进一步：这事现在完全是我和\Person{Octavius}之间的事。如果他愿意让我这个持有那种见解的人与他辩论，他马上就会明白，在同伴中争论比以哲学家的方式激烈交锋更容易。让我们坐在那些为保护浴场而堆砌、延伸入海的石堤上，这样我们既能休息旅途，又能更专注地辩论。}听了他的话，我们坐了下来，他们从两侧掩护我，使我处在三人中间。这并不是出于礼节、等级或荣誉，因为友谊总是一视同仁；而是让我作为仲裁者，靠近双方以便倾听，并坐在中间将两人分开。于是\Person{Cæcilius}这样开始：——\par

% structure: p0010 source=h2 target=section reason=relative-heading-rank; base=h2

\section{第五章 论题：\Person{Cæcilius}首先提醒众人，在人事中一切事物都是可疑和不确定的，因此应当惋惜的是，基督徒——他们大多未受训练且目不识丁——竟敢于对万物主宰和天主威严作任何确定的论断；由此他论证世界不受任何天主眷顾治理，并得出结论，最好恪守传承的宗教形式。}

尽管对你而言，\Person{马可}兄弟，我们所特别探究的主题并无疑问，因为你已仔细了解两种生活，拒绝了其一而赞同了另一，但在此情形下，你必须调整心态，犹如一位最公正的法官把握平衡，不可偏袒任何一方（多于另一方），免得你的裁决似乎并非源于我们的论证，而出自你个人的感知。因此，你若以一新进且对双方皆无知者的身份审判我，则不难表明，人事中的一切皆可疑、不确定、未定，万物皆属于概然而非真实。因此不足为奇，有些人因对彻底探究真理感到厌倦，便草率屈服于某种意见，而不愿以不懈的勤奋坚持探寻。故而所有人必然愤慨，所有人必然痛心，竟有一些人——既未受过学问训练，又对文学无知，甚至不请低贱技艺——敢于就万物的本性及（天主的）尊威妄下断论，而历代如此众多的教派对此仍存疑惑，连哲学本身也仍在思量。此非无理；因为人类智力的平庸远不能胜任（探究）神圣之事的能力，以致我们既不得知晓，也不许探究，且不可亵渎地强求那悬于上天之上的事物，或深深潜于地下的事物；并且，按照古代智者的那句神谕，我们若能深刻认识自己，便似乎可称得足够幸福、足够谨慎。但即便我们沉溺于无谓而徒劳的劳作，逾越谦卑所当有的界限，并且，虽倾向尘世，却以大胆的野心超越上天本身和诸星，至少让我们不要以虚妄可怕的见解纠缠这一错误。纵使万物的种子当初是由一种自身结合万物的自然凝聚而成——此处哪位天主是创造主？纵使世界的各部分是由偶然的相遇联合、分解、塑造而成——哪位天主是设计者？纵使火点亮了星辰；纵使（其物质的轻盈）悬起了穹苍；纵使其物质以重量奠定了大地；纵使海水自湿气而来，那么，这宗教从何而来？这恐惧从何而来？这迷信为何？人，以及一切出生、赋予生命、被滋养的动物，不过是元素的自愿凝结，而人及一切动物又复归其中，被分解、消散。如此，万物复返其本源，回归自身，并无任何工匠、审判者或创造者。如此，火种聚集，造成其他的太阳，又其他太阳，永远照耀。如此，大地的蒸气蒸腾，造成雾气不断增长，雾气凝聚汇集，造成云层升高；云层下降时，造成降雨、刮风、冰雹哗哗坠落；或当云层相撞时，造成雷声咆哮，闪电发红，雷霆闪耀。因此它们随处落下，冲向山峰，击中树木；毫无选择地摧毁神圣与世俗之地；击打恶人，也常常击打虔诚之人。我何必提那多变而不定的风暴，其中对万物的攻击毫无秩序与区分地翻腾？——在沉船中，善人与恶人的命运混杂，他们的报应混淆？——在火灾中，无辜者与有罪者一同毁灭？——当天空的瘟疫玷污一地区时，所有人无区别地死去？——当战争的热浪肆虐时，往往是更优秀的人倒下？在和平时期也一样，邪恶不仅与善人（的境遇）等量齐观，甚至备受尊崇，以致对许多人而言，你不知是该最憎恶他们的堕落，还是该羡慕他们的幸福。但若世界是由天主眷顾和某位神祇的权柄所治理，\Person{法拉里斯}和\Person{狄奥尼修斯}就绝不会配得为王，\Person{鲁提里乌斯}和\Person{卡米卢斯}就绝不会该遭放逐，\Person{苏格拉底}就绝不会该服毒。看那结果子的树，看那已发白的庄稼，已垂下的葡萄，被雨水摧毁，被冰雹击打。因此，要么一个不确定的真理向我们隐藏并被扣留；要么，更可信的是，在这些多变而任性的偶然中，不受法律约束的命运，正统治着我们。\par

% structure: p0012 source=h2 target=section reason=relative-heading-rank; base=h2

\section{论点：所有民族，尤其是罗马人，崇拜其神祇的目标，是为了对其崇拜获得对整个大地的至高统治权。}

既然命运是确定的，或本性是不确定的，那么，作为真理的大司祭，接受祖先的教导，遵守传授给你们的宗教，敬拜你们最初从父母那里学会畏惧而非熟悉认识的神祇，不是更虔敬、更好吗？不要对神明妄加论断，而要相信你们的祖先，他们在世界本身诞生的年代尚处未受教之时，就已经配得上拥有神祇，或为保佑他们，或为治理他们。因此，我们看到在所有帝国、行省和城市中，每个民族都有自己的国家敬拜仪式，敬拜他们本地的神祇：如厄琉息斯人敬拜色勒斯；弗里吉亚人敬拜玛特尔；厄庇道洛斯人敬拜埃斯库拉庇乌斯；加色丁人敬拜贝鲁斯；叙利亚人敬拜阿斯塔特；陶里斯人敬拜狄安娜；高卢人敬拜墨丘利；罗马人敬拜所有神明。因此，他们的权力和权威占据了整个世界的圆周；因此，他们将帝国扩展到了太阳的轨道和海洋的边界之外；因为他们在战争中实践着宗教的勇毅；因为他们用圣礼的宗教、贞洁的处女、众多的尊荣和司铎的名号来巩固他们的城市；因为当被围困和夺取时，除了卡皮托利山几乎全部沦陷，他们仍敬拜那些其他民族在愤怒时会鄙视的神祇；而穿过高卢人的战线，高卢人对他们迷信的胆大感到惊奇，他们不带武器，却以宗教的崇拜为武装；在敌人的城市被夺取、仍处胜利狂怒之时，他们崇敬被征服的神祇；他们四处寻找外邦人的神祇，并使之成为自己的；他们甚至为未知的神祇和玛内斯建立祭台。因此，由于他们承认所有民族的宗教制度，他们也配得了统治权。因此，他们崇敬的永恒进程得以延续，不因漫长时光的流逝而削弱，反而增强，因为古时常赋予仪式和庙宇与岁月同等的神圣性。\par

% structure: p0014 source=h2 target=section reason=relative-heading-rank; base=h2

\section{七、论点：罗马的征兆和占卜被忽视带来恶果，被遵守带来好运。}

我们的祖先也并非偶然（因为我敢暂时推定\Emph{辩论之点}，如此安全地犯错）通过遵守占兆、咨询内脏（占卜）、设立圣礼或奉献庙宇而在事业上取得成功。思考一下书籍的记载。你立刻会发现，他们开启了各种宗教的仪式，要么是为了回报神祇的恩宠，要么是为了避免灾难性的愤怒，要么是为了平息那已经膨胀和爆发的愤怒。见证伊代阿母亲，她的到来既证明了主妇的贞洁，又将城市从敌人的恐惧中解救出来。见证那对骑士兄弟的雕像，它们正如在湖上显现的那样被奉献，他们骑着气喘吁吁、口吐白沫和冒烟的马，在得胜的同一天宣布了战胜波斯人的消息。见证因一个平民的梦而重新举行的被冒犯的朱庇特的赛会。德基乌斯家族的奉献也是一个公认的见证。库尔提乌斯也是如此，他用他的骑士身份的重担或荣耀填满了深渊的开口。此外，比我们更频繁的是，占兆被轻视时，就见证了神祇的临在：因此阿利亚河是一个不吉利的名字；因此克劳狄乌斯和尤尼乌斯的战役不是对抗迦太基人的战役，而是一场致命的沉船。因此，为了使特拉西美努斯湖因罗马人的血而肿胀变色，弗拉米尼乌斯轻视了占兆；并且为了让我们再次从帕提亚人那里要回我们的军旗，克拉苏既应得又嘲笑可怕姐妹的诅咒。我省略了许多古老的故事，并略过诗人们关于神祇的出生、恩赐和奖赏的诗歌。此外，我匆匆掠过神谕所预言的命运，免得古事在你们看来过于荒诞。看看罗马城赖以保护和武装的神庙和神祇的神龛：它们因居住其中的神祇——这些神祇常驻于其中——而更加庄严，而非因崇拜的徽章和奉献而富丽。因此，先知们充满神，与神结合，预先收集未来，为危险提供警告，为疾病提供药物，为受苦者提供希望，为不幸者提供帮助，为灾难提供安慰，为劳苦提供缓解。即使在我们的休息中，我们也能看到、听到、认识到那些我们在白天不敬地否认、拒绝和背弃的神祇。\par

% structure: p0016 source=h2 target=section reason=relative-heading-rank; base=h2

\section{八、论点：狄奥多鲁斯、狄亚哥拉斯和普罗塔哥拉斯的亵渎鲁莽决不可被默许，他们要么想完全摆脱神祇的宗教，要么至少削弱它。但更不可容忍的是那些躲藏和避光的基督徒，他们拒绝神祇，并且害怕死后死亡，却不怕现在就死。}

因此，既然万民关于不朽诸神存在的共识已经确立，尽管其本性或起源仍不确定，我绝不容忍任何人因某种我不知为何的狂妄和不信神的智慧而自我膨胀，企图动摇或削弱这如此古老、如此有益、如此健全的宗教，即使此人可能是\Place{昔兰尼}的\Person{狄奥多罗斯}，或在他之前的\Place{墨利安}的\Person{狄亚戈拉斯}——古人曾赋予他们\OriginalTerm{无神论者}{Atheist}的称号；这两人都断言没有神，从而夺去了统治人类的一切敬畏和所有的尊崇；然而，他们决不会以他们所假称的哲学之名和权威，在这不虔敬的教义上得逞。当雅典人驱逐了\Place{阿布德拉}的\Person{普罗泰戈拉}，并当众焚烧他的著作时，只因他关于神性的论述是故意而非亵渎性的，那么，为什么就不能痛惜这样的事：一些堕落的、不法的、绝望的派别（请容忍我相当自由地运用我所承担的论辩的力量）竟向诸神发怒？他们从最底层渣滓中聚集了最无知的人，以及轻信而因其性情的软弱而屈从的妇女，组成一个渎神阴谋的群体，这群体由夜间集会、严肃斋戒和非人的食物联结在一起——不是通过任何神圣礼仪，而是通过需要补赎的行为——一群躲藏避光、在公共场合沉默、在角落里饶舌的人。他们藐视庙宇如同墓地，拒绝诸神，嘲笑圣物；可怜的他们，若被允许，则同情祭司；他们半裸着身体，藐视尊荣和紫袍。哦，奇异的愚昧和难以置信的狂妄！他们藐视当下的痛苦，尽管他们害怕那些不确定的、将来的痛苦；而当他们害怕死后死去时，却不怕现在就死：如此，一种欺骗性的希望以复活的慰藉安抚了他们的恐惧。\par

% structure: p0018 source=h2 target=section reason=relative-heading-rank; base=h2

\section{第九章 论点：基督徒的宗教是愚蠢的，因为他们崇拜一个被钉死的人，甚至崇拜他受刑的刑具本身。据说他们还崇拜驴的头，甚至他们父的本性。他们通过屠杀和婴孩的血得以入教，在不知羞耻的黑暗中不分彼此地混杂在一起。}

如今，随着更邪恶的事日益增多，堕落的习俗日复一日地蔓延，那亵渎集会的可憎场所正在全世界成熟。无疑，这种联盟应当被根除和咒骂。他们通过秘密标记和徽章彼此识别，几乎在认识之前就彼此相爱。到处都有某种淫欲的宗教混杂其中，他们不分彼此地互称兄弟姐妹，以便那并不罕见的放荡行为因这神圣名称的介入而成为乱伦：正是如此，他们虚妄而愚昧的迷信在罪中夸耀。若非背后有真相，明智的报告也不会谈及如此重大繁多的、需要以辩护作为前言的事情。我听说他们崇拜驴的头——那最卑贱的受造物，被某种我不知道的愚昧信念所祝圣——这是对这种习俗的相称而合适的宗教。有人说他们崇拜他们祭司长和祭司的\OriginalTerm{生殖器}{virilia}，并仿佛崇拜他们共同父母的本性。我不知道这些事是否虚假；当然，怀疑适用于秘密和夜间的仪式；而有人将他们的仪式解释为指向一个因邪恶而受极端痛苦惩罚的人，以及指向十字架那致命的木头，这为堕落和邪恶的人提供了合适的祭坛，好让他们崇拜配得之物。至于初学弟子入教的故事，既应被憎恶，又众所周知。一个裹着面粉的婴孩，以欺骗不知情者，被放在即将被他们的仪式玷污的人面前：这个婴孩被那位年轻弟子杀死——他仿佛被鼓励对面粉表面进行无害的击打，却用隐秘而隐蔽的伤口施以杀害。他们贪婪地——哦，恐怖！——啜饮其血；他们急切地分其肢体。他们因这祭品而彼此约束；凭借这作恶的意识，他们誓守相互沉默。这样神圣的仪式比任何亵渎更可憎。至于他们的宴席，众所周知，到处都有人谈论；甚至我们\OriginalTerm{基尔塔人}{Cirtensian}的演说也见证了这一点。在某个庄严的日子，他们带着所有孩子、姐妹、母亲、各种性别和年龄的人，聚集在宴会中。在那里，盛宴过后，当团体变得热烈，乱伦情欲的狂热因醉酒而升温时，一只被拴在吊灯上的狗被人用抛出绳索长度之外的一点碎肉激怒而猛冲跳跃；于是那有意识的灯被推翻，在无耻的黑暗中熄灭，可憎情欲的结合使他们陷入命运的未知。虽然并非所有人在事实上都是乱伦的，但在意识上所有人都同样乱伦，因为每个人的欲望都寻求在个体行为中可能发生的一切。\par

% structure: p0020 source=h2 target=section reason=relative-heading-rank; base=h2

\section{第十章 论点：无论基督徒崇拜什么，他们都竭力隐瞒：他们没有祭台，没有庙宇，没有公认的偶像。他们的神，如同犹太人的神，据说是一位；他们既不能看见也不能显示他，却认为他是恶意的、不安分的、不合时宜地好管闲事的。}

我有意略过许多事，因为所提的已经太多；这一切或其中大部分为真，可由其卑劣宗教的隐晦性推知。因为为何他们如此费力地隐藏并掩盖他们崇拜的一切？尊贵的事总是乐于公开，而罪恶则秘而不宣。他们为何没有祭台、没有圣殿、没有公认的雕像？为何从不公开谈论、从不自由聚集？除非他们所崇拜并隐藏的，要么是应受惩罚的，要么是可耻的。再者，那\Emph{独一}、孤寂、被遗弃的天主，从何处而来，是谁，又在哪里？没有一个自由的民族、没有一个王国、甚至罗马的迷信也不曾认识祂。孤独而悲惨的犹太民族崇拜一位天主，一位专属于他们的天主；但他们公开崇拜祂，有圣殿、祭台、牺牲和礼仪；然而这位天主竟如此无力无能，连同祂自己的特殊民族一起，被罗马的神祇奴役。但基督徒又虚构何等奇事、何等怪物！——他们说他们那位既不能显示也不能看见的天主，仔细查考所有人的品性、行为，甚至言语和隐秘思想；祂到处奔走，处处临在；他们将祂描绘成麻烦、不安，甚至无耻地好奇，因为祂临在于所行的一切事，往来于一切地方，尽管祂忙于整体，却无法关注细节，也不可能在忙于细节时照料整体。什么！因为他们威胁要用烈火毁灭整个世界，乃至宇宙及其所有星辰，难道他们是在图谋毁灭万物？——好像由自然神律所建立的永恒秩序会被扰乱，万有的联盟会被打破，天体的结构会瓦解，那包含并维系万有的架构会倾覆似的。\par

% structure: p0022 source=h2 target=section reason=relative-heading-rank; base=h2

\section{第十一章　论题：除主张未来世界将遭大火焚毁外，他们还承诺我们身体的复活：义人将享有最幸福的生命，不义者将受极刑。}

他们不满足于这种狂悖之见，还附加上老妪的虚构故事：他们说死后、骨灰与尘土中将要复活；不知凭何自信，他们轮流相信彼此的谎言——你会以为他们已经复活了。这是双重的祸患和双倍的疯狂：一方面宣告天穹和星辰的毁灭，而我们则听其自然；另一方面却向已死已灭的我们——既生即灭的我们——许诺永生！无疑正因如此他们诅咒我们的火化堆，谴责火葬，仿佛任何身体即使不被火焰吞噬，也会因年岁而最终化为尘土，仿佛不论身体被野兽撕裂、被海洋吞噬、被泥土覆盖、还是被火焰卷走，都毫无区别；因为对尸体而言，若尚有感觉，每一种埋葬方式都是惩罚；若无感觉，则在毁灭的迅速中反得解脱。被此谬误所骗，他们向自己——作为善人——应许死后福乐永生的生命，而向他人——作为不义者——应许永罚。有许多事本可补充，若非我的论述催促我结束。我已证明他们自己是不义的，且不须再费心证明；尽管即便我承认他们是义人，你们的共识也与多数人的意见一致，认为罪与无辜归于命运。因为我们所做的一切，正如一些人归于命运，你们则归于天主：如此，按你们的宗派，人并非出于自愿，而是被拣选而愿意。因此你们想象出一个不义的审判者，祂惩罚人并非因其意愿，而是因其命运。但我很想知道，你们复活时是否带着身体？若带，是带着同样的身体还是更新的身体？不带身体？那么，据我所知，就没有心智、灵魂、生命了。带着同样的身体？但这身体先前早已毁坏。带着另一个身体？那便是新人的出生，而非原人的复原；然而如此漫长的岁月、无数的时代已经流逝，有哪一个个体从死者中返回过？——哪怕是像\Person{普罗忒西劳斯}那样获准停留片刻，好让我们以此为榜样而相信？所有这类不健康信仰的虚构和空洞的安慰，被说谎的诗人在甜美诗句中戏弄，都被你们可耻地重述，你们毫不怀疑地相信你们的天主。\par

% structure: p0024 source=h2 target=section reason=relative-heading-rank; base=h2

\section{第十二章　论题：况且，基督徒自己死后将遭遇何事，可从他们现今已一无所有、备受最沉重灾祸与苦难的事实预见。}

至少你们没有从当前事物中取经验，看出那徒劳期望的空幻许诺如何欺骗你们。可怜的人们，从你们的命运，趁着你们还活着，想想死后什么在威胁你们。看，你们的一部分——按你们所说，是更大更好的一部分——正匮乏、寒冷、辛勤劳作和饥饿；天主却容忍此事，他假装；他要么不愿意要么不能帮助他的子民；因此他或是软弱或不公正。你们这些梦想着死后不朽永生的人，当被危险摇撼，被热病消耗，被痛苦撕裂时，难道不感受到你们的真实状况吗？难道不承认你们的软弱吗？可怜的人，你们是否不情愿地被说服了你们的软弱，却不肯承认？但我忽略那些人人相同的事。看，对你们有威胁、刑罚、折磨和十字架；这些不再是崇拜的对象，而是要受的折磨；还有火，你们既预言又恐惧。那位在你们复活时能帮助你们的天主在哪里？既然他在今生不能帮助你们。难道罗马人，没有你们天主的任何帮助，不是照样治理、统治、享受整个世界，并辖制你们吗？而你们在这期间，在悬疑焦虑之中，克制自己不去体面的享乐。你们不去观看表演，不参与公共演出，拒绝公共宴席，憎恶神圣竞赛，回避祭台上已尝过的肉和奠过的酒。因此你们畏惧你们所否认的神祇。你们不以花环装饰头颅，不以香料妆点身体，把油膏保留给葬礼，甚至拒绝给你们的坟墓献花环——苍白、颤抖的存在，甚至值得我们神祇的怜悯！因此，你们如此可怜，既不会复活，也不能在当下生活。所以，如果你们有一点智慧或谦逊，就停止窥探天空的领域、世界的命运和秘密吧：留意脚下就足够了，尤其是对未受教育、粗野、乡下的平民而言：那些没有能力理解世俗事务的人，更被剥夺了讨论神圣事物的能力。\par

% structure: p0026 source=h2 target=section reason=relative-heading-rank; base=h2

\section{第十三章 论辩：\Person{凯基利乌斯}最终得出结论，新宗教应被摒弃；且我们不可轻率判断悬疑之事。}

\Quote{不过，如果你们想要从事哲学，你们中任何一个足够伟大的人，若能效法智慧之王\Person{苏格拉底}，便应如此。那个人对于天上之事被问及时的回答是众所周知的：\InnerQuote{\Emph{我们之上的事情与我们无关}。}因此，他确实从神谕那配得独特智慧的见证，这神谕他自己也预感到，他被置于众人之上，并非因为他发现了一切，而是因为他认识到自己一无所知。因此，承认无知是智慧的高峰。由此，\Person{阿尔克西拉乌斯}、后来很久的\Person{卡尔内亚德}以及许多学园派在最重要的问题上产生了安全的怀疑，在这种哲学流派中，无学者可以谨慎地做到很多，而有学者则可以荣耀地做到。什么！\Person{西蒙尼德斯}这位抒情诗人的犹豫不决难道不值得所有人钦佩和效仿吗？当\Person{希耶罗}僭主问他以为神是什么、像什么时，他先要求一天来思考；然后将答复推迟了两天；然后，在催逼下，又加了一天；最终，当僭主询问如此长时间延误的原因时，他回答，他的研究时间越长，真理对他变得越模糊。依我之见，不确定的事情应当保持原样。在如此众多伟大的人物还在思考之时，我们不应轻率鲁莽地提出相反的观点，以免引入幼稚的迷信，或推翻一切宗教。}\par

% structure: p0028 source=h2 target=omitted reason=duplicate-page-title

到此为止，\Person{凯齐留}说了这一番话；他愉快地微笑着（因为冗长演说的激烈已缓解了他愤怒的热忱），补充道：\Quote{\Person{屋大维}，身为\Person{普劳图斯}的族裔，他曾是磨坊主中的首领，却仍是哲学家中最卑微者，他竟敢对此作何答复？}\Quote{克制一下，}我说，\Quote{不要对他自鸣得意；因为在双方更充分地辩论这主题之前，你不配因你演说的和谐而得意，尤其是你的推理追求的是真理而非赞美。尽管你的演说以其微妙的变化令我愉悦，但我深受触动，并非关乎当前的讨论，而是关乎整个辩论的方式——即真理的状况大多会随讨论的力量而改变，甚至清晰的雄辩能力也是如此。这众所周知是由于听众的轻信所致：他们被言辞的诱惑从对事物的注意中分心，对所说的一切毫无区别地赞同，不能分辨虚假与真理；他们不知道，即使在看似不可信的事物中也常有真理，而在貌似真实中也有虚假。因此，他们越是相信大胆的主张，就越常被更聪明的人说服，从而不断被自己的轻率所欺骗。他们将评判者的过错转嫁于不确定的抱怨；以致于他们谴责一切，宁可让所有事悬而不决，也不愿对可疑之事做出决断。因此我们务必小心，不要让我们因此招致对所有言论的憎恨，正如许多较为单纯的人被引导去诅咒和憎恶一般人。因为那些轻信的人被他们本以为是值得信赖的人所欺骗；之后，由于相似的错误，他们开始怀疑每个人都是邪恶的，甚至惧怕那些他们本应视为卓越的人。所以现在我们忧虑——因为每件事都可能有两面之词，一方面真理大多模糊，另一方面有一种奇妙的巧辩，有时以其言语的丰富模仿公认证明的确信——因此我们尽可能仔细地衡量每个细节，以便我们准备赞赏敏锐的同时，却能选择、认可并采纳正确的事物。}\par

% structure: p0030 source=h2 target=omitted reason=duplicate-page-title

\Quote{你正在退出，}\Person{凯齐留}说，\Quote{一位宗教评判者的职责；因为你插入一个非常重要的论点来削弱我诉状的力量，这是非常不公平的，因为\Person{屋大维}面前有我说过的每件事，完整无损，只要他能反驳它们。}\Quote{你所指责的，}我说，\Quote{除非我错了，我是为了共同的利益而提出来的，以便通过细致的审查，我们不是根据辞藻华丽的风格，而是根据事情本身的坚实性来权衡我们的决定。我们不应再因你的抱怨而被分散注意力，而是要以绝对的沉默聆听我们朋友\Person{雅努阿里乌}的答复，他正在向我们招手。}\par

% structure: p0032 source=h2 target=omitted reason=duplicate-page-title

\Person{奥克塔维乌斯}于是这样开始说：我定会尽我所能、竭我之力来说话，而你们也必须与我一同努力，用真实的言辞之流冲淡非常刺耳的互相指责之声。我一开始也不掩饰，我的朋友\Person{纳塔利斯}的意见摇摆不定，那样游移、含糊、滑溜，以至于我们不得不怀疑，是你的信息混乱了，还是它纯因错误而前后摇摆。因为他时而相信诸神，时而又对此犹豫不决；因此我答复的直接目的由于他主张的不确定性而变得更加不确定。但至于我的朋友纳塔利斯——我不允许，也不相信有任何诡辩——他的淳朴远非狡诈的伎俩。那么又怎样呢？正如一个不知正路的人，恰巧一条路分岔成许多条，因为不知路，他就焦虑不安，既不敢选择某一条路，也不敢尝试所有路；同样，一个人若对真理没有坚定不移的判断，就如他不信的怀疑四散，他怀疑的意见也不确定。因此，\Person{凯基利乌斯}同样被潮水抛来抛去，在互相矛盾抵触的事物间颠簸，就不足为奇了；但为了使这不再发生，我要揭发并反驳所有已说的话，无论它们多么多样化，只确认并赞许真理；今后他再不可怀疑或游移。既然我的兄弟爆发出了这样的言语，说他悲伤，他烦恼，他愤慨，他遗憾没有学问、贫穷、不熟练的人竟讨论神圣之事；让他知道，所有人都是同等地被生，具有推理和感知的能力和才能，没有年龄、性别或地位的偏颇。他们也不是靠运气获得智慧，而是由自然植入；此外，那些哲学家本身，或其他作为技艺发现者而声名显赫的人，在凭心智技能获得显赫名声之前，也被视为平民、未受教化、半裸。因此，富人依附于他们的财富，习惯于更多地凝视他们的金子而不是上天，而我们这类人，虽然贫穷，却既发现了智慧，又将教导传授给别人；由此可见，才智不是赐予财富的，也不是通过学习获得的，而是与心灵的形成一同产生的。因此，如果有人探究、思考、表达关于神圣之事的想法，没有什么可愤怒或悲伤的；因为所需要的不是争论者的权威，而是争论本身的真理；言论越是缺乏技巧，推理就越明显，因为它没有被浮夸的雄辩和优雅所渲染；而实际上，它是靠正义的准则来支撑的。\par

% structure: p0034 source=h2 target=section reason=relative-heading-rank; base=h2

\section{第十七章 论证：人确实应当认识自己，但除非他首先认识万物的整个范围以及\Person{天主}本身，否则无法获得这种认识。并且，从世界本身的构造和陈设来看，每个赋有理性的人都认为世界是由\Person{天主}建立的，并由祂治理和管理。}

我也不拒绝承认\Person{凯基利乌斯}所竭力主张的主要事项之一，即人应当认识自己，环顾四周，看清自己是什么、从何而来、为何存在；无论是由元素聚合而成，还是由原子和谐地形成，抑或更确切地说，是由天主创造、形成并赋予生命。而正是这一点，若非通过对宇宙的探究，我们便无法寻求和查考；因为万物如此连贯、如此紧密相连、相辅相成，以至于除非你勤勉地考察天主的本性，否则你必然对人性一无所知。你也不能很好地履行你的社会责任，除非你认识那共同享有的世界社群，尤其是因为在这方面我们与野兽不同：它们俯首向地，生来只知关注食物；而我们，面容昂然，目光仰望上天，正如我们的交谈和理性，借此我们认识、感知并效法天主，我们既没有权利也没有理由对那形成于我们眼前和感官中的天上荣耀无知。因为在低处寻找本应在高处寻找的东西，无异于最严重的亵渎。因此，那些否认整个世界这一构造是由神圣理智所完善，并断言它是由某些偶然粘合在一起的碎片堆积而成的人，在我看来，既没有心智，也没有感觉，甚至没有视力。因为当你举目望天，察看下方和周围的事物时，还有什么比这更明显、更公认、更显而易见呢：即有一位拥有最卓越理智的神明，整个自然被他所激发、所推动、所滋养、所治理？请看天空本身，它扩展得如此广阔，旋转得如此迅速，无论是在夜间被星宿点缀，还是在白昼被太阳照亮，你立刻就会知道至高治理者的奇妙而神圣的平衡在其中运作。再看一年，它是如何由太阳的运转形成的；再看一月，月亮如何在它的盈、亏、暗中驱动它运转。我该说什么关于黑暗与光明的反复交替呢？我们如何因此获得了劳作与休息的交替恢复？诚然，关于星宿的更冗长的论述必须留给天文学家，无论是关于它们如何引导航行的路线，还是带来耕种或收获的季节，这些事物中的每一个不仅需要一位至高工匠和完美的理智，不仅需要创造、构建和安排，而且它们本身若非通过最高的理智和理性就无法被感知、察觉和理解。怎么！当季节和收获的秩序以稳定的变化被区分时，它岂不为其创造者与父亲作证？春天带着它的花朵，夏天带着它的收获，秋天带着它令人感激的成熟，以及冬天的橄榄采摘，都是必要的；而这一秩序很容易被打乱，除非它由最高的理智所建立。现在，需要多么伟大的天主眷顾，以免只有冬天用霜冻摧残，或只有夏天用炎热灼烧，而是介入秋天和春天的适中温度，使回归的岁月那不可见且无害的过渡悄然滑过！请仔细看大海；它被海岸的法则所束缚。无论何处有树木，请看它们如何从大地的深处获得活力！考虑海洋；它以交替的潮汐涨落。看那泉源，它们如何奔流不息！凝视河流；它们总是沿规整的河道滚滚向前。我为何要谈论那适当地排列的山峰、丘陵的斜坡、平原的广阔？我为何要谈论有生命的受造物彼此对抗的多种保护？——有的武装以角，有的以牙齿作篱，以爪为鞋，以刺为钩，或凭借足部的速度获得自由，或凭借翅膀的飞翔能力？我们自身形态的美丽尤其承认天主是它的工匠：我们直立的体态，我们仰望的面容，我们的眼睛仿佛置于高处作为眺望，以及我们所有其他感官仿佛布置在城堡中。\par

% structure: p0036 source=h2 target=section reason=relative-heading-rank; base=h2

\section{论点：而且，天主不仅照顾整个世界，也照顾其各个部分。藉由独一真神的命令，万事万物得以治理，这由地上王国的例证得以证明。但是，尽管祂无限而广大——祂何其伟大，只有祂自己知道——既不能被我们看见，也不能被我们称呼，然而祂的荣耀最清楚可见，正是在我们放下所有称号之时。}

若要逐一列举实例，说来话长。人身上没有哪一个肢体不是为了需要和美观而设计的；更奇妙的是，所有肢体形式相同，但每个都有某些线条的微调，于是我们每个人彼此相异，却又大体相似。我们为何出生？生育的欲望意味着什么？难道不是天主所赐，使乳房在胎儿成熟时充满乳汁，使柔弱的幼体借着丰沛乳水的滋养而成长吗？天主不仅照顾整个宇宙，也照顾它的各个部分。\Place{不列颠}缺乏阳光，却被环绕的海水之温暖所滋润；\Place{尼罗河}调和了\Place{埃及}的干旱；\Place{幼发拉底河}灌溉了\Place{美索不达米亚}；\Place{印度河}弥补了雨水的不足，据说它同时播种和浇灌了东方。现在，如果你进入一所房屋，看到一切精致、有序、装饰华丽，你必定会相信有一位主人治理着它，而且他本人比所有这些美好的事物更美好。同样，在这世界的房屋中，当你观察天空和大地、它的眷顾、它的秩序、它的法则时，请相信有一位宇宙的主宰和父亲，他远比星辰和世界的各部分更荣耀。除非，也许——既然对天主眷顾的存在没有疑问——你认为这是一个探究的问题：天上的王国是由一个权力统治，还是由多个统治；而这件事本身并不难阐明，对于一个思考尘世帝国的人来说，尘世帝国的实例无疑是从天上来的。何时有过王权的同盟是始于诚信而不以流血告终的？我不提\Place{波斯人}，他们从马的嘶鸣中收集统治的预兆；我也不引述那个绝对死亡的\Person{忒拜兄弟}的寓言。关于双胞胎（\Person{罗慕路斯}和\Person{勒莫斯}）的故事，涉及牧羊人和小屋的统治，是众所周知的。女婿与岳父的战争遍布全世界；如此伟大帝国的命运容不下两个统治者。看看其他事物：蜜蜂有一个蜂王；羊群有一个领头的；牛群中有一个统治者。你能相信在天上至高的权力有分争，那个真实而神圣的帝国的全部权威被分裂吗？因为显然，万有的父、天主既无开始也无终结——他生出万有，却使自己永恒——他在世界之前，他自己就是世界。他借着言语安排一切，以智慧布置，以能力成全。他不能被看见——他比光更明亮；不能被把握——他比触摸更纯净；不能被估量——他超越一切感觉；无限、无量，他的伟大只有他自己知道。但我们的心太有限，无法理解他，因此我们恰当地估量他，当我说他超出估量时。我要说出我的感受。凡以为知道天主的伟大的人，是在缩小它；凡不愿缩小它的人，就不认识它。你也不可向天主求名。天主就是他的名。当众多需要通过名称的特殊特征来区分个体时，我们需要名字；但对于唯一的天主，天主这个名字就是全部。如果我称他为父，你会判断他为属世的；如果称他为王，你会怀疑他是属肉欲的；如果称他为主，你定会认为他是可朽的。除去名字的附加，你将看见他的荣耀。怎么！难道在这件事上我没有得到所有人的同意吗？我听见百姓向天举手时，只喊出：\Quote{哦，天主}，\Quote{天主伟大}，\Quote{天主真诚}，\Quote{若天主准许}。这是百姓自然的言谈，还是一个承认基督信仰者的祈祷？那些以朱庇特为首的人，虽然在名字上错了，但在权力的统一上却是一致的。\par

\EditorialNote{"马嘶声"指\Person{大流士·希斯塔斯佩斯}的事例。凡读过希罗多德著作的人都记得这个故事，这无疑是历史上最非凡的故事之一——一匹马竟选出了一位伟大的君主，而此人的生命对人类命运产生了不小的影响。一个狡猾的马夫的确是这个选举的策划者，正如在这些事件中，历史的秘密弹簧常被隐藏；但如果这个故事不完全是寓言，那么天上打雷的巧合是最值得注意的。它似乎预示着天主眷顾的支配，以及天主将人的愚昧（不亚于愤怒）转化为对天主的赞美的能力。参见希罗多德《历史》第三卷第八十六章。}\par

% structure: p0039 source=h2 target=section reason=relative-heading-rank; base=h2

\section{第19章 论据：此外，诗人们称他为神与人的父亲、万物的创造者、他们的心灵与精神。而且，甚至更优秀的哲学家们也几乎得出与基督徒相同的关于天主统一的结论。}

我听到诗人们也宣称\Quote{诸神和人的一位父}，并且凡人的理智如此，正如万有之父所指定的一日。\Person{曼图亚的马罗}怎么说？难道不是更清楚、更贴切、更真实吗？他说：\Quote{起初，内在的精神滋养着，灌注的理智搅动着天和地}，以及世界的其他部分，\Quote{由此产生了人和牛群}，以及其他各种动物。同一诗人在另一处将那理智和精神称为天主。因为他的话是：\Quote{因为那位天主遍及所有陆地、海洋的区域和深邃的天际；人类和牲畜来自祂；雨和火也来自祂。}此外，天主被我们宣告为不是别的，而是理智、理性和精神吗？如果乐意，让我们回顾一下哲学家的教导。尽管他们的论述方式多样，但你会发现他们在这件事上一致同意。我略过那些未经训练的古人们，他们因其话语而配被称为智者。让米利都的泰勒斯为首，因为他首先讨论天界之事。同一位米利都的泰勒斯说水是万物的本原，但天主是从水中形成万物的那理智。啊！对水和精神的描述，比人类所曾发现的更高尚、更尊贵。这是由天主启示给他的。你看，这位原始哲学家的意见与我们的完全一致。随后阿那克西美尼，以及阿波罗尼亚的第欧根尼，断定无限而不可测度的空气即是天主。他们在神性上的认同也与我们相似。而阿那克萨哥拉的描述则是，天主据说是无限理智的运动；毕达哥拉斯的天主则是那遍及万物的普遍本性、来回穿梭且专注的灵魂，所有动物的生命也源自祂。众所周知，色诺芬尼宣称天主是拥有一颗理智的无限整体；安提斯泰尼则说，民众多神，但有一位自然的天主是万物的首领；克塞乌西普斯承认一种自然的有生命的力量为天主，万物由它治理。德谟克利特怎么说？尽管他是原子的最初发现者，但他不也特别把作为形式基础和理智的自然称为天主吗？斯特拉托本人也说天主就是自然。此外，伊壁鸠鲁，这个虚构无所事事的神明或根本无神的人，却仍然将自然置于万有之上。亚里士多德则变化不定，但仍赋予一种单一力量：他时而说理智是天主，时而说世界是天主，时而又将天主置于世界之上。本都的赫拉克利德斯也以各种方式将神圣理智归于天主。泰奥弗拉斯托斯、芝诺、克吕西波和克莱安特斯本人确实有多种意见，但都归结为同一个事实：天主眷顾的统一性。因为克莱安特斯论述天主时，有时当作理智，有时当作灵魂，有时当作空气，但大多时候当作理性。他的老师芝诺认为自然和天主的法则，有时是空气，有时是理性，为万物的本原。此外，通过将朱诺解释为空气，朱庇特为天，尼普顿为海，伏尔甘为火，并以同样方式指明普通民众的其他神祇为元素，他有力地驳斥并克服了公众的错误。克吕西波几乎说同样的话。他相信一种神圣的力量、理性的自然，有时是世界和必然的命运，就是天主；他效法芝诺，对赫西俄德、荷马和俄耳甫斯的诗歌进行生理学解释。此外，巴比伦的第欧根尼的教导是：阐述并论证朱庇特的诞生、密涅瓦的起源等等，都是对其他事物的名称，而非诸神。因为苏格拉底的门徒色诺芬说，真天主的形状不可见，因此不应寻求。斯多葛派的阿里斯托说，祂完全不能被理解。两人都感受到了天主的威严，但都绝望于理解祂。柏拉图对天主的论述更为清晰，无论在事情本身还是他所用的名称上；他的论述本是全然超性的，若不是偶尔被掺入世俗信仰所玷污。因此，在他的 \WorkTitle{蒂迈欧篇} 中，柏拉图的天主按其名字就是世界的父、灵魂的创造者、天上和地上事物的制造者；他宣称，因祂那超乎寻常、令人难以置信的能力，发现祂是困难的，而一旦发现祂，在公开场合谈论祂是不可能的。我们的意见也几乎相同。因为我们知道并谈论一位万有之父的天主，除非被问及，绝不在公开场合谈论祂。\par

% structure: p0041 source=h2 target=section reason=relative-heading-rank; base=h2

\section{第二十章 论点：但如果世界由天主眷顾治理并统摄于独一天主旨意之下，愚昧的敌意不该将我们卷入与它一同错误的迷途；即使它迷恋自身的寓言，已引入了荒谬的传统。同样清楚显明，对诸神的崇拜始终愚妄不虔，因为古昔之人尊崇其君王、名将和技艺发明者，因其赫赫功业而敬之如神。}

我已经陈述了几乎所有哲学家的意见，他们更为卓越的荣耀在于指出了有一位天主，尽管有众多名称；以致任何人或许会认为，要么基督徒如今是哲学家，要么哲学家那时已经是基督徒。但如果世界受天主眷顾治理，并由唯一天主的旨意引导，那么无知古人的意见，纵然因其自身的寓言而欢喜迷醉，也不该将我们引入相互认同的错误，因为这种意见已被其自身的哲学家的意见所反驳，而这些哲学家既有理性的权威，又有古人的权威支持。因为我们的祖先对谎言抱有如此轻易的信仰，以至于他们轻率地相信甚至其他怪物是奇妙的奇迹；一个多形的\OriginalTerm{斯库拉}{Scylla}，一个多形态的\OriginalTerm{喀迈拉}{Chimæra}，一个从其吉祥的伤口中再次升起的\OriginalTerm{许德拉}{Hydra}，以及\OriginalTerm{半人马}{Centaurs}，马与骑手缠绕在一起；凡是传闻允许捏造的任何东西，他们都完全愿意听信。我何必提及那些老妇的寓言，说人从人变成鸟兽，又从人变成树木和花草？这些事情，如果曾经发生，就会再次发生；而因为它们现在不可能发生，所以从未发生过。同样，对于众神也是如此，我们的祖先漫不经心、轻信地、以未经训练的单纯来相信；他们虔敬地崇拜自己的国王，渴望在死后以可见的形式见到他们，急于在雕像中保存他们的记忆，那些仅仅作为安慰而采用的东西变得神圣。于是，在世界尚未因商业而开放，各民族尚未混淆其礼仪和习俗之前，每一个特殊的民族都尊崇其创始人、或杰出的领袖、或比其性别更勇敢的谦逊女王、或任何才能或技艺的发现者，作为值得纪念的公民；如此，对逝者给予奖赏，对后来者树立榜样。\par

% structure: p0043 source=h2 target=omitted reason=duplicate-page-title

请读\OriginalTerm{斯多葛学派}{Stoics}的著作，或智者的作品，你将与我一同承认这些事实。由于他们德行的功绩或某些恩赐，\Person{欧赫美尔}断言他们被尊为神；他列举了他们的生日、国家、墓冢，并在各个省份指出狄克忒的\Person{朱庇特}、德尔斐的\Person{阿波罗}、法罗斯的\Person{伊西斯}和厄琉息斯的\Person{刻瑞斯}这些事迹。\Person{普罗迪科斯}谈到一些人被擢升到诸神之中，因为他们在漂泊中发现新种类的产物，有助于人的需要。\Person{珀耳塞斯}也作同样的哲理论述；他并补充说，被发现的果实和这些果实的发现者被以相同的名称称呼；正如喜剧作家的段落所说：\Quote{维纳斯没有\Person{巴克斯}和\Person{刻瑞斯}就会僵冷。}著名的马其顿人\Person{亚历山大大帝}，在他写给母亲的一篇著名文献中写道，由于畏惧他的权势，司祭向他泄露了诸神原本是人的秘密：他告诉母亲，\Person{武尔坎}是万物的始祖，然后是\Person{朱庇特}的族系。而你们见到\Person{伊西斯}的燕子和铙钹，以及你们的\Person{塞拉皮斯}或\Person{奥西里斯}的空坟，其肢体四散。那么再来看圣礼本身，以及它们真正的奥秘：你们会看到悲惨的死亡、不幸、葬礼，以及可怜诸神的哀伤与恸哭。\Person{伊西斯}与她的\Person{犬头神}和秃头司祭一同哀悼、哭泣、寻找她失去的儿子；可怜的伊西斯崇拜者捶打胸膛，模仿这最不幸母亲的悲伤。不久，当小男孩被找到时，\Person{伊西斯}欢喜，司祭们雀跃，\Person{犬头神}发现者夸耀，而他们年复一年地要么失去找到的，要么找到失去的。为所崇拜的悲伤，或为所悲伤的崇拜，岂非可笑？然而这些本是埃及的礼仪，如今却成了罗马的。\Person{刻瑞斯}手持点燃的火炬，被一条蛇环绕，焦虑而忧心地追踪被偷走的\Person{珀耳塞福涅}的足迹——她在游荡中被劫，成为败坏者。这就是厄琉息斯秘仪。\Person{朱庇特}的圣礼又如何？他的乳母是只母山羊，他作为婴儿被从贪婪的父亲那里带走，以免被吞噬；\Person{科律班忒斯}的铙钹敲出铿锵的喧闹，以免父亲听到婴儿的哭泣。\Place{丁狄摩斯}的\Person{库柏勒}——我羞于提及——她无法引诱她那通奸的情人（不幸的是，他对她欢喜）行淫，因为她自己，作为众神之母，又丑又老，便阉割了他，无疑是为了使那阉人成为神。因着这个故事，\Person{伽利}也通过阉割自己身体的方式崇拜她。当然，这些不是圣礼，而是酷刑。诸神的形貌又如何？难道不证明你们神祇的可鄙与可耻本色吗？\Person{武尔坎}是跛足之神，瘸腿；\Person{阿波罗}在多少世代后仍无胡须；\Person{埃斯库拉庇乌斯}胡须浓密，尽管他是永远年轻的\Person{阿波罗}之子；\Person{尼普顿}有海绿色的眼睛；\Person{密涅瓦}有灰蓝色的眼睛；\Person{朱诺}有牛眼；\Person{墨丘利}有翼足；\Person{潘}有蹄足；\Person{萨图尔努斯}脚戴镣铐；\Person{雅努斯}有两张面孔，仿佛可以回头看路；\Person{狄安娜}有时是女猎手，裙子高高束起；作为\Place{以弗所}的，她有很多丰乳；作为夸张的\Person{特里维亚}，她有三头多手，可怖。你们的\Person{朱庇特}本身又如何？他时而雕像为无胡须，时而树立为有胡须；被称为\Person{哈蒙}时，他有角；称为\Person{卡皮托利努斯}时，他挥舞雷电；称为\Person{拉提亚里斯}时，他溅满鲜血；称为\Person{费雷特里乌斯}时，无人接近；且不提众多\Person{朱庇特}，\Person{朱庇特}的怪异形貌与他的名号一样多。\Person{厄里戈涅}被吊在绞索上，以便作为童贞女在群星间闪耀。\Person{卡斯托耳们}轮流死亡，以便活着。\Person{埃斯库拉庇乌斯}为使升为神，被雷电击中。\Person{赫丘利}为脱去人性，被\Place{厄塔山}的火焰焚尽。\par

% structure: p0045 source=h2 target=section reason=relative-heading-rank; base=h2

\section{第二十二章 论点：此外，这些起初由无知之人所编造的寓言，后来被他人，主要是诗人所颂扬，他们的权威对真理造成了不小的损害。通过这类虚构以及更具吸引力的谎言，年轻人的心灵被败坏，因而他们可悲地终老于这些信念，尽管另一方面，只要他们愿意寻求，真理对他们而言是显而易见的。}

这些寓言和错误，我们从无知的父母那里学来，而更严重的是，我们在学习和教导中——尤其是在诗人的诗句中——将它们精心加工，诗人们以其权威尽可能地损害了真理。正因如此，\Person{柏拉图}在他所建立的理想国中，公正地将那位他曾赞美并加冕的著名诗人\Person{荷马}驱逐出去。因为正是在\Place{特洛伊}战争中，他尤其容许你们的神祇（尽管他嘲弄他们）仍干预人类的事务和行为：他将他们聚在一起争斗；他伤了维纳斯；他捆绑、击伤并赶走了玛尔斯。他叙述说，朱庇特被布里亚柔斯释放，以免被其他神祇捆绑；他用血雨哀悼自己的儿子萨耳珀冬，因为他无法将他从死亡中夺回；并且，受维纳斯腰带的诱惑，他与妻子朱诺同寝，比他与那些奸情的恋人更加急切。在其他地方，赫拉克勒斯清除粪肥，而阿波罗为阿德墨托斯放牧牲畜。然而，尼普顿为拉俄墨冬筑墙，不幸的建造者没有得到他工作的报酬。然后，朱庇特的霹雳在铁砧上与埃涅阿斯的兵器一同锻造，尽管在朱庇特生于\Place{克里特}之前很久，就已经有天、霹雳和闪电；独眼巨人既无法模仿，朱庇特自己也害怕真正霹雳的火焰。我何必提起玛尔斯和维纳斯被发现的奸情，以及朱庇特对伽倪墨得斯的暴行——\Quote{如你们所说}，这是一件在天上被尊崇的恶行？所有这些都是为了给人类的恶行赋予某种权威。通过这些虚构的故事以及其他类似的故事，以及更迷人的谎言，孩童的心灵被败坏；随着同样的传说附着于他们，他们甚至在成年之后仍然如此；这些可怜的人，尽管真理显而易见——只要他们愿意寻求——却在同样的信仰中老去。因为所有古代的作家，无论希腊人还是罗马人，都已说明，\Person{萨图尔努斯}——这个种族和众人的始祖——是一个人。\Person{内波斯}知道这点，\Person{卡西乌斯}在他的历史中也知道；\Person{塔卢斯}和\Person{狄奥多罗斯}也这样说。这个\Person{萨图尔努斯}被从\Place{克里特}驱赶，因害怕他那愤怒的儿子，来到\Place{意大利}，受到\Person{雅努斯}的款待，教导那些无知而粗野的人许多事情——因为他多少有些希腊人的教养和文雅——例如，传授文字、铸造钱币、制造工具。因此，他更愿意将他藏身的地方——因为他曾安全地隐藏（latent）在那里——称为\Place{拉丁姆}；并以自己的名字命名一座城市为\Place{萨图尔尼亚}，\Person{雅努斯}则命名为\Place{贾尼科洛}，这样他们各自将自己的名字留给后人的记忆。因此，那确实是一个逃亡的人，确实是一个隐藏的人，一个人的父亲，由人所生。然而，他被宣布为大地或天的儿子，因为在\Place{意大利}人中，他的父母不明；正如直到今日，我们称那些突然出现、不知来自何处的人为从天上降下来的，那些出身微贱、无名的人为大地之子。他的儿子朱庇特在父亲被驱逐后统治\Place{克里特}。他在那里死去，在那里有儿子。直到今日，朱庇特的洞穴仍被参观，他的坟墓也被展示，而通过他那些神圣的仪式本身，他被证明是人。\par

% structure: p0047 source=h2 target=omitted reason=duplicate-page-title

无须逐一细述每一个案，也无须展开那一族类的全部系列，因为在其初祖身上，他们的必死性已被证实，并且按照他们传承的法则，必定会流溢到其余的人身上——除非你们幻想他们死后成了神；正如因\Person{普洛库卢斯}的伪誓，\Person{罗慕路斯}成了神；因\Place{毛里塔尼亚人}的善意，\Person{尤巴}成了神；还有其他的君王也被奉为神明，他们被祝圣，不是基于他们神性的信仰，而是为了尊崇他们曾行使的权力。而且，这个名号被加给那些不愿承受它的人。他们渴望保持人的状态，害怕被奉为神；尽管他们已经年老，却并不愿意。因此，神既不能由死人造成，因为神不能死；也不能由受生之人造成，因为凡受生的都会死。但那既无升起也无沉落者才是神圣的。因为，如果他们是受生的，为何如今不再出生呢？——除非，或许\Person{尤庇特}已经年老，\Person{朱诺}已不能生育，\Person{弥涅耳瓦}在生育之前已经白头。或者，这种生殖过程已经停止，因为不再有人相信这类传说？此外，如果神能创造，他们就不能朽坏：我们会有比所有人加起来还多的神；以致现今，天不能容纳他们，气不能接受他们，地不能承载他们。由此明显可见，那些我们读到受生、又知道死去的人，原本就是人。因此，谁会怀疑平民向这些人的被祝圣的像祈祷并公开崇拜呢？因为无知者的信念和心智被艺术的完美所欺骗，被黄金的光辉所蒙蔽，被白银的光泽和象牙的洁白所眩惑？但是，如果有人用心想一想每一尊像是用什么工具和机械制成的，他就会羞愧自己竟畏惧被工匠随意加工以制造神明的物质。因为木神，或许是一堆柴或一根不吉利的木材，被悬挂、被砍伐、被劈开、被刨平；铜神或银神，往往来自不洁的器皿，正如\Place{埃及}国王所做的那样，被熔炼、被锤打、在铁砧上锻造；石神被砍凿、雕刻、被某个堕落的人打磨，它在诞生时感觉不到所受的伤害，后来也感觉不到来自你们崇敬的崇拜；除非那石头、木料或银子还不是神。那么，神何时开始存在呢？看哪，它被熔化、被打造、被雕刻——还不是神；看哪，它被焊接、被组装——被竖立起来，仍然不是神；看哪，它被装饰、被祝圣、被祈祷——终于成了神，当人选择它如此，并为此目的奉献了它。\par

% structure: p0049 source=h2 target=section reason=relative-heading-rank; base=h2

\section{第二十四章 论点：他简要地指出，在庆祝某些神的秘仪时，人们遵循了何等可笑、猥亵和残忍的仪式。}

无言的动物对你们的神的判断是多么更真实啊！老鼠、燕子、鸢都知道它们没有感觉：它们啃咬它们，践踏它们，坐在它们上面；除非你们赶走它们，它们就在你们神的嘴里筑巢。蜘蛛甚至在他的脸上织网，从他的头上垂丝。你们擦拭、清洗、刮磨，并保护和畏惧你们自己制造的东西；而你们中没有一人认为在崇拜天主之前应当先认识天主；他们不加思索地想要服从祖先，宁愿成为他人错误的附庸，也不信任自己；因为他们对自己所畏惧的一无所知。因此贪婪在金子和银子中被祝圣了；因此虚空的雕像的形状被确立了；因此\Place{罗马人}的迷信兴起了。如果你们重新思考这些神的仪式，有多少事是可笑的，又有多少事是可悲的！赤裸的人们在严冬中奔跑；有些人戴着帽子行走，带着旧盾牌，或打鼓，或领着他们的神沿街乞讨。有些神殿允许一年进入一次，有些则完全禁止入内。有一个地方男人不得进入，有些地方对妇女神圣不可侵犯：奴隶参加某些仪式也是需要赎罪的罪行。有些圣地由有一个丈夫的妇女加冕，有些则由有多个丈夫的妇女加冕；而能够数出最多通奸次数的人，反被最虔诚地追求。什么！一个用自己的血奠酒、以自己的伤口祈求（他的神）的人，如果完全世俗，难道不比以这种方式虔诚更好吗？而一个被割去羞处的人，以这种方式试图平息天主的愤怒，对天主是何等的冒犯！因为，如果天主想要阉人，祂可以如此创造他们，而不会事后使他们如此。谁看不出心神不健全、理解力薄弱而堕落的人在这些事上是愚蠢的，而且犯错的人如此众多，彼此提供了相互的庇护？在这里，普遍疯狂的辩护就是疯狂人数众多。\par

% structure: p0051 source=h2 target=section reason=relative-heading-rank; base=h2

\section{第二十五章 论点：他接着指出，\Person{凯基利乌斯}错误地断言\Place{罗马人}通过适当遵守此类迷信而获得了对整个世界的统治权。相反，\Place{罗马人}在其起源上是由罪行聚集而成，并以其残酷的恐怖而壮大。因此，\Place{罗马人}之所以如此伟大，并非因为他们虔诚，而是因为他们亵渎神明而未受惩罚。}

然而，你或许会说，正是那种迷信本身为罗马人赋予、增长并巩固了他们的帝国，因为他们取胜与其说是靠勇气，不如说是靠宗教与虔诚。无疑，罗马人那光辉而正义的国度，自其成长中的帝国之初便已开始。难道他们起先不是聚集在一起，以罪恶为堡垒，凭自身的凶残之威而壮大的吗？最初的人群如同聚集到一个避难所一般。被遗弃的人、放荡者、乱伦者、凶手、叛徒，蜂拥而至；为了使他们的统帅与统治者\Person{罗慕路斯}本人比他的子民更胜一筹，他犯下了弑亲罪。这就是那宗教国度的最初征兆！不久，他们抢夺、玷污并毁灭了异邦的处女，那些已经许配、注定为妻的少女，甚至还有一些刚结婚的年轻女子——这是前所未闻的——然后与她们的父母，即她们的岳父们开战，流了亲属的血。还有什么比这更不虔诚、更胆大妄为、更安全无虞的呢？因为对罪行的自信本身就是安全的。如今，驱逐邻人离开土地，推翻邻近的城市及其庙宇和祭坛，将他们掳为俘虏，以他人的损失和自身的罪恶而壮大，这是\Person{罗慕路斯}之后其余国王和后来领袖们共同接受的训练方式。因此，罗马人所拥有、耕种、占有的一切，都是他们胆大妄为的掠夺品。他们所有的庙宇都是用暴力的战利品建造的，即用城市的废墟、神祇的战利品和司铎的屠杀而建。这就是侮辱和蔑视，是对被征服的宗教的屈服，是在战胜它们之后又崇拜它们作为俘虏。因为崇拜你凭借武力夺得的东西，就是奉圣渎为神圣，而不是奉神祇。因此，罗马人每凯旋一次，就受玷污一次；他们从列国获得的战利品越多，从神祇那里掠取的掳物也就越多。所以，罗马人之所以如此伟大，并非因为他们虔诚，而是因为他们毫无惩罚地亵渎神明。因为在战争中，他们根本无法得到那些他们与之兵戎相见的诸神的帮助；他们开始崇拜那些被战胜的神祇，即他们先前所挑战过的。但这些神祇对罗马人有什么帮助呢？他们尚且不能保护自己的敬拜者免遭罗马人的武力。因为我们知道罗马本地的神——\Person{罗慕路斯}、\Person{皮库斯}、\Person{提贝里努斯}、\Person{孔苏斯}、\Person{皮卢姆努斯}和\Person{皮库姆努斯}。\Person{塔提乌斯}发现了\Person{克洛基娜}并崇拜她；\Person{霍斯提利乌斯}发现了敬畏与苍白之神。后来，我不知道谁奉献了热病神：这就是滋养那座城市的迷信——疾病和健康不良的状态。的确，\Person{阿卡·劳伦提娅}和\Person{弗罗拉}，臭名昭著的妓女，也必须算作罗马人的疾病和诸神之一。无疑，正是这样的神祇扩展了罗马人的统治，与那些被列国崇拜的神祇相对抗：因为针对他们自己的子民，色雷斯的\Person{玛尔斯}、克里特的\Person{朱庇特}、\Person{朱诺}（时而\Place{阿尔戈斯}，时而\Place{萨摩斯}，时而\Place{迦太基}）、\Place{陶里斯}的\Person{狄安娜}、\Person{伊达山圣母}，以及那些埃及的——不是神祇，而是怪物——都没有帮助过他们；除非或许在罗马人中，贞女的贞洁更高贵，或司铎的宗教更神圣：尽管事实上在非常多的贞女中，不贞洁受到了惩罚——她们无疑在\Person{维斯塔}不知情的情况下，与男子过于轻率地交往；至于其余的，她们的免罚并非由于她们贞洁的更好保护，而是由于她们不贞的更好运气。而且，哪里还有比在祭台和神龛之间司铎安排通奸更好的地方呢？哪里更有拉皮条的讨论，或更多的暴力密谋呢？最后，燃烧的情欲在圣殿看守的小室里得到满足，比在妓院本身还要频繁。而且，早在罗马人之前，按天主的安排，亚述人、米底人、波斯人、希腊人以及埃及人曾统治过，尽管他们没有任何大司祭、\OriginalTerm{阿尔瓦兄弟}{Arvales}、\OriginalTerm{萨利司铎}{Salii}、\OriginalTerm{维斯塔贞女}{Vestals}、\OriginalTerm{占卜师}{Augurs}，也没有关在笼子里的鸡，靠它们的啄食或禁食来治理国家最高事务。\par

% structure: p0053 source=h2 target=omitted reason=duplicate-page-title

现在我要论及你们煞费苦心收集的那些罗马占兆和占卜，你们为之作证，说它们被忽视时带来恶果，被遵守时带来好运。诚然，\Person{克洛狄乌斯}、\Person{弗拉米尼乌斯}和\Person{尤尼乌斯}因此失去了他们的军队，因为他们认为不值得等待由贪食鸡的啄食所给出的那非常庄严的预兆。但\Person{雷古卢斯}又如何？他难道不是遵守了占卜却被俘了吗？\Person{曼奇努斯}持守了他的宗教义务，却被置于轭下，被交出去。\Person{保卢斯}在\Place{坎尼}也有贪食鸡，却与共和国的大部分一起覆灭了。\Person{盖乌斯·凯撒}蔑视那些反对他在冬季前航行去\Place{阿非利加}的占卜和占兆，于是反而更轻易地航行了并征服了。但我还要继续讲多少关于神谕的事？死后，\Person{安菲阿拉俄斯}回答了关于未来的事，尽管他活着的时候并不知道自己会因为一只手镯被妻子出卖。盲人\Person{提瑞西阿斯}看见了未来，虽然他看不见现在。\Person{恩尼乌斯}编造了皮堤亚的\Person{阿波罗}关于\Person{皮洛士}的回答，尽管阿波罗早已停止作诗；而他那个谨慎而模棱两可的神谕，恰好在人们开始变得更有教养且更不轻信的时候失效了。而\Person{德摩斯梯尼}，因为他知道那些回答是伪造的，便抱怨皮媞亚在\Emph{亲腓力化}。但有时，确实，甚至占兆或神谕也曾触及真相。虽然在众多虚谎之中，偶然可能显得仿佛在模仿先见；然而我将接近那错误和悖逆的真正源头，一切黑暗皆从那里涌出，并更深地挖掘它，更显明地揭露它。有一些伪诈而游荡的灵体，被尘世的污点和情欲从它们属天的活力中贬低。如今这些灵体，因被罪恶压倒和沉浸而失去了本性的纯朴，作为其灾祸的慰藉，它们自身已灭亡，却不停息地使别人灭亡；它们自身已败坏，却将败坏的错误注入别人；它们自身已与天主疏离，却通过引入低劣的迷信使别人与天主分离。诗人们知道那些灵体是魔鬼；哲学家们论述它们；\Person{苏格拉底}知道这事，他根据身边一个魔鬼的点头和决定，要么拒绝要么承担事务。贤士们也不仅知道存在魔鬼，而且，无论他们假装行什么奇迹，都是借魔鬼之手而行；借着他们的呼求和交流，他们展示奇妙的把戏，使不存在的事物显现，或使存在的事物隐没。在这些魔术师中，无论从口才还是从事迹上都是第一的\Person{索斯提尼斯}，不仅以合宜的威严描述了真天主，还描述了那些作为天主的仆役和信使的天使，甚至是真天主本身。他知道这加深了对祂的敬畏，即因对它们的主的点头和注视的畏惧而战栗。同一个人也宣称魔鬼是属地的、游荡的、与人类为敌的。\Person{柏拉图}说了什么？他相信发现天主是一件困难的事。难道他不也毫不犹豫地谈论天使和魔鬼吗？在他的\WorkTitle{会饮篇}中，他难道不也努力解释魔鬼的本性吗？因为他认为那是一种介于可朽与不朽之间的实体——即介于身体与精神之间的中介，由属地的沉重和属天的轻盈混合而成；由此他也告诫我们爱的欲望，并说它被塑造并滑入人的胸膛，搅动感官，塑造情感，注入情欲的热火。\par

% structure: p0055 source=h2 target=section reason=relative-heading-rank; base=h2

\section{第27章 论题：重述。毫无疑问，这里有一个错误的源头：魔鬼潜伏在雕像和图像之下，他们出没于神庙，他们激活内脏的纤维，引导鸟的飞行，掌管签，倾泻包含虚假回应的神谕。这些事并非来自天主；但当他们以真天主之名被命令时，他们被迫承认，并从附身的身体中被驱逐。因此他们赶紧逃离基督徒的附近，并在那些在认识他们之前就开始恨他们的外邦人心中激起对他们的仇恨。}

这些不洁之灵，即魔鬼，正如贤士、哲学家及柏拉图所示，隐藏在雕像和图像之下，凭借其气息获得如同在场神祇的权威；同时，它们被吸入先知体内，居住于神庙，有时鼓动内脏的纤维，控制鸟的飞行，指挥签卜，是充满许多谎言之神谕的根源。因为它们既受骗，又骗人；既不知简单的真理，又为自毁而不承认所知。因此它们将人从天上拉下，使人离开真天主而转向物质事物；扰乱生活，使所有人不得安宁；也秘密潜入人体，以其精灵的微妙，假装疾病，惊扰心神，扭曲肢体；迫使人们崇拜它们，饱享祭台的香烟或牲畜的牺牲，好让它们通过解除自己所束缚的，似乎治愈了病症。这些疯狂的狂人，你看见他们在公共场所狂奔，他们自身也是没有圣殿的先知；他们如此狂怒，如此癫狂，如此旋转。在他们里面也有同样的魔鬼煽动，但他们疯狂的原因不同。同样原因也产生了你不久前所说的事情：朱庇特在梦中要求恢复他的竞技，卡斯托尔骑马出现，一只小船跟随主妇的腰带指引。许多人，甚至你们自己的一些人，都知道魔鬼自己承认关于自身的一切，当我们用言辞的折磨和祈祷之火将它们从身体中驱赶出来时，它们常常如此。萨图尔努斯本人，以及塞拉皮斯、朱庇特，和你们崇拜的任何魔鬼，因痛苦而说出它们是什么；它们确实不会说谎而有损自己，尤其当你们中有人在场时。既然它们自己是魔鬼的见证，相信它们对自身的真实告白吧；因为当被唯一真天主所驱逐时，这些可怜的东西不情愿地在身体中战栗，要么立即跳出，要么逐渐消失，视病患的信德或治愈者的恩宠所助。因此它们从近处的基督徒逃遁，而远距离时它们却通过你们在集会中骚扰他们。这样，它们被引入无知者的心灵，暗中由于恐惧播下对我们的仇恨。因为人自然恨他所怕的，如果可能，伤害他怕过的人。因此它们占据心灵，阻塞心扉，使人们在不认识我们之前就开始恨我们；以免一旦被认识，他们要么效法我们，要么不能定我们的罪。\par

% structure: p0057 source=h2 target=section reason=relative-heading-rank; base=h2

\section{论点：他们不仅煽动对基督徒的仇恨，还指控他们犯有可怕的罪行，至今无人证实。这是魔鬼的作为。因为魔鬼既散布又传播虚假报告。基督徒被错误地指控亵渎、乱伦、奸淫、杀亲；此外，确实且真实的是，外邦人自己实际上也犯了同样的或类似甚至更大的罪行。}

然而，你们这样对未知和未经审查的事妄下判断，是多么不公正！当你们悔改时，请相信我们这些过来人，因为我们从前也和你们一样，曾是盲目愚钝的，也认为相同的事：即基督徒崇拜怪物、吞食婴儿、参与乱伦的宴席。我们并未察觉这类谣言总是由那些好事者散布，从未被调查或证实；在如此长的时间里，竟无一人出来揭发，以获得对其罪行的宽恕，甚至为揭露罪行而赢得青睐；而且，基督徒被指控时既不脸红也不害怕，他只后悔从前不是基督徒，由此可见这些事并非如此邪恶。然而，当我们自告奋勇辩护和庇护某些亵神乱伦者，甚至弑亲者时，却认为这些基督徒根本不值得听。有时，当我们假装可怜他们时，反而对他们更加残酷，以至在他们认罪时折磨他们，好让他们否认，免得灭亡；我们使用颠倒的审讯，不是为了查明真相，而是为了强迫他们说谎。如果有人因软弱而受不住痛苦，屈服并否认自己是基督徒，我们便对他表示好感，好像他通过简单否认这个名号就立刻赎清了一切罪行。难道你们不承认我们曾像你们现在这样感受和行事吗？若是由理性而非恶魔的煽动来判断，本应强迫他们不否认自己是基督徒，而应承认自己犯有乱伦、可憎之事、亵渎圣仪、杀害婴儿等罪。因为那些恶魔就是用这些类似的故事，在无知者面前填满对我们的恐惧，使他们厌恶害怕。这并不奇怪，因为人的流言飞语总是靠散布谎言来滋养，一旦真相大白，它便消散。这就是恶魔的作为，因为它们既播种又培育了谣言。由此产生了你们所说的听闻：驴头在我们中被视为神圣之物。谁会愚蠢到崇拜它呢？谁又会更愚蠢到相信它被崇拜？除非你们竟在厩中连同你们的埃波娜一起供奉整头驴，并虔诚地与伊西斯一同吞食这些驴。你们还献祭并崇拜牛头和公羊头，并奉祀混合了山羊和人的神，以及狗脸和狮脸的神。你们岂不也与埃及人一起崇拜和饲养阿比斯牛吗？你们并不谴责为尊崇蛇、鳄鱼、其他野兽、飞鸟和鱼类而设立的圣礼，若有人杀了这些神中的任何一个，甚至会被处死。这些埃及人，连同你们中的许多人，畏惧伊西斯还不如惧怕洋葱的辛辣，敬畏塞拉比斯还不如避开身体秽物所发出的最卑贱声音。那些编造我们崇拜祭司器官的谣言的人，其实是将他自己所行的事转嫁给我们。（因为这些可能是他们淫秽的圣事，他们那里所有性别以所有器官出卖，他们那里所有淫秽被称为风雅；他们嫉妒娼妓的自由，他们舔舐中间男子，淫荡的嘴附着在私处，这些人口舌恶劣，即使沉默也令人生厌，他们对自己的淫秽厌倦先于羞愧。）可憎！他们使自己遭受如此邪恶的行为，没有哪个时代如此柔弱能忍受，也没有哪种奴役如此残酷能被迫承受。\par

% structure: p0059 source=h2 target=section reason=relative-heading-rank; base=h2

\section{第二十九章 论点：基督徒崇拜因罪行被钉在十字架上的人，这更不真实，因为他们不仅相信他是无辜的，而且合理地相信他是天主。但另一方面，异教徒祈求被自己奉为神明的国王的神力；他们向偶像祈祷，恳求他们的守护神。}

这些以及诸如此类的无耻之事，我们甚至不能听；再用更多言语为自己辩护这些指控是可耻的。因为你们假想那些贞洁谦逊的人做了我们除非你们证明你们自己也确有其事，否则根本不信会有人做的事。你们将崇拜罪犯和他的十字架归咎于我们的宗教，这与真理相去甚远，竟以为一个罪犯配得上，或一个尘世之人能被相信是天主。那将全部希望寄托于凡人者实在可怜，因为人的一切帮助都随人消亡而终结。埃及人确实为自己选一个人来崇拜；他们只向他求情；他们事事请教他；向他献祭；他对于别人是神，对于自己却无论愿意与否确实是人，因为他即使欺骗别人，也无法欺骗自己的良心。此外，虚假的谄媚无耻地奉承君主国王，不是当作伟大杰出的男子（这样才正当），而是当作神；然而，对显赫之人致以敬意更为真实，对至善之人表示爱戴更为愉快。因此他们呼唤他们的神祇，祈求他们的偶像，恳求他们的守护神，即他们的魔灵；对着朱庇特的守护神起假誓，比对着国王的守护神起假誓更安全。至于十字架，我们既不崇拜也不渴望。你们这些供奉木制神像的人，或许将木制的十字架当作你们神祇的一部分来崇拜。因为你们的军旗、旌旗和营旗，除了是镀金装饰的十字架外，还能是什么？你们的胜利纪念品不仅模仿简单十字架的形状，还模仿被钉在上面的人。我们当然在船帆鼓起航行时，在船桨展开滑行时，在抬起战车轭时，自然见到十字架的标记；当人以纯洁的心灵伸开双手敬拜天主时，也是如此。所以，十字架的标记或因自然理性而存在，或你们自己的宗教也与之相关。\par

% structure: p0061 source=h2 target=section reason=relative-heading-rank; base=h2

\section{第30章 辩论：关于基督徒杀害婴儿喝其血的故事是赤裸裸的诽谤。但外邦人既残酷地遗弃他们刚出生的孩子，又用残忍的堕胎在出生前就毁灭他们。基督徒既不允许看见也不允许听到杀人。}

现在，我愿会见那个说或相信我们通过杀戮婴儿和饮其血而加入教会的人。你能相信如此柔软、微小的身体能够承受那些致命伤吗？任何人能流淌、倾倒并吸干一个幼小者和一个几乎尚未存在的人的鲜血吗？除了敢于去做的人，无人能相信。我看到你们有时将你们所生的孩子遗弃给野兽和飞鸟；另一些时候，你们用可悲的死亡方式勒死他们。有些妇女通过饮用药物制剂，在她们的腹中扼杀未来人的源头，从而在生育前就犯下弑亲之罪。这些事无疑来自你们神祇的教导。因为\Person{撒杜恩}没有遗弃自己的孩子，而是吞噬了他们。有理由在\Place{非洲}的一些地区，父母将婴儿献祭给他，用爱抚和亲吻压抑他们的哭声，以免献祭一个哭泣的牺牲。此外，在\Place{彭托斯}的\Person{陶里人}中，以及向\Place{埃及}的\Place{布西里斯}，献祭宾客是一种神圣的仪式，而\Person{高卢人}向\Person{墨丘利}献上人的，或者说非人的祭品。\Place{罗马}的献祭者活埋了一个\Place{希腊}男子和一个\Place{希腊}女子，一个\Place{高卢}男子和一个\Place{高卢}女子；直到今天，\Person{拉提亚里斯·朱庇特}仍被他们以谋杀来崇拜；而且，作为\Person{撒杜恩}之子实至名归，他充满了邪恶和罪人的血。我相信他本人教导\Person{喀提林}以血的契约谋反，教导\Person{贝罗娜}用一口人血浸润她的神圣仪式，并教导人们用一个人的血来治疗癫痫，即以更糟糕的疾病来治疗。他们也不亚于那些从竞技场吞噬野兽的人，浑身沾满血迹，或者以人的肢体或内脏养肥自己。对我们来说，看见或听见杀人都是不合法的；我们如此避免人血，以至于我们连可食用动物的血也不用于我们的食物。\par

% structure: p0063 source=h2 target=section reason=relative-heading-rank; base=h2

\section{第31章 辩论：关于我们的宴会被乱伦玷污的指控完全违背一切可能性，而明显的是，外邦人实际上犯有乱伦罪。基督徒的宴席不仅谦逊，而且节制。事实上，乱伦的欲望如此闻所未闻，以至于对许多人来说，即使是两性之间谦逊的交往也会引起脸红。}

至于乱伦的宴席，恶魔的阴谋虚构了一个巨大的谎言来反对我们，以污秽我们谦逊的光荣，通过一种恶劣的丑闻所引起的憎恶，使得人们在查究真相之前，就被一种可憎的指控的恐怖而远离我们。在这方面，你们自己的\Person{弗朗托}正是这样做的：他并没有作为提出指控的人提供证据，而是作为一个修辞学家散布指责。因为这些事实际上源于你们自己的民族。在\Place{波斯}人中，儿子与母亲之间的混杂交往是被允许的。在\Place{埃及}人和\Place{雅典}人中，与姐妹结婚是合法的。你们的记录和你们的悲剧，你们愉快地阅读和聆听，以乱伦为荣：同样，你们崇拜乱伦的神祇，他们与母亲、女儿、姐妹交合。因此，有理由在你们当中频繁发现乱伦，并且不断被允许。可怜的人们，你们甚至可能在不自觉的情况下冲进非法之事：因为你们混杂地散布你们的欲望，因为你们到处生育孩子，因为你们常常甚至将在家出生的孩子遗弃给他人，不可避免的是，你们必定回到你们自己的孩子那里，并迷失于你们自己的后代。这样，你们继续着乱伦的故事，即使你们没有意识到自己的罪行。但是，我们保持谦虚不是在外表上，而是在心里，我们乐意遵守单一婚姻的约束；在生育的愿望中，我们知道要么有一个妻子，要么一个都没有。我们实行共享宴席，这些宴席不仅谦逊，而且节制：因为我们不沉溺于娱乐，也不以酒延长我们的盛宴；但我们以严肃、贞洁的谈论、甚至更贞洁的身体（我们中的许多人是未受玷污的）来节制我们的欢乐，享受而非夸耀身体永久的童贞。事实上，他们离沉迷于乱伦欲望如此之远，以至于对有些人来说，即使是（想到）两性之间谦逊的交往也会引起脸红。我们拒绝你们的荣誉和紫袍，也不立即站在最低层民众的水平上；如果我们都识别出一种善，我们并不挑剔，而是以我们个人生活的同样宁静聚集在一起；我们不在角落里喋喋不休，尽管你们在公开场合听到我们要么脸红，要么害怕。而且，我们的人数日复一日地增加，这不是错误指控的根据，而是值得称赞的见证；因为，在一种公平的生活方式中，我们的实际人数持续存在，不减反增，陌生人增加了它。简而言之，我们不像你们所假设的那样，用某种小小的身体标记来区分我们的人，而是用无辜和谦逊的标志很容易地加以区分。因此，我们彼此相爱，令你们遗憾，以一种相互的爱，因为我们不知道如何恨。因此，我们彼此称呼，令你们嫉妒，为弟兄：作为由同一天主和父母所生的人，信仰中的同伴，希望中的共同继承人。然而，你们不认识彼此，在彼此的仇恨中残忍；也不承认彼此为弟兄，除非是为了杀害弟兄。\par

% structure: p0065 source=h2 target=section reason=relative-heading-rank; base=h2

\section{第32章 辩论：也不能说基督徒隐藏他们所崇拜的，因为他们没有圣殿和祭台，因为他们确信天主不能被任何圣殿所限制，也不能制作任何祂的肖像。但祂无处不在，看见一切，甚至我们心中最隐秘的思想；我们生活在祂近旁，在祂的保护之下。}

但你认为，如果我们没有圣殿和祭台，我们就隐藏了我们所敬拜的对象吗？然而，若你思想正确，人本身就是天主的肖像，我还能造什么天主的像呢？祂的工程所塑造的整个世界都不能容纳祂，我该为祂建造什么圣殿呢？而我这个人四处居住，难道要把如此伟大威严的力量关闭在一座小小的建筑里吗？祂岂不更应被供奉在我们的理智中，被祝圣在我们的内心深处吗？我该向主献上牺牲和祭品，如同祂为我的用处所生产的，让我把祂自己的恩赐还给祂吗？当适合祭献的牺牲是良好的心性、纯洁的理智和真诚的判断时，这是忘恩负义的。因此，培养纯洁的人祈求天主；实践正义的人向天主奉献；远离欺诈行为的人平息天主的义怒；从危险中拯救人的人献上了最蒙悦纳的牺牲。这就是我们的祭献，这就是我们敬拜天主的礼仪；因此，在我们中间，最正义的人就是最虔诚的人。然而，我们所敬拜的天主，我们既不展示，也不看见。实在说来，正因如此我们相信祂是天主：我们能够意识到祂，却不能看见祂；因为在祂的工程中，在世界的一切运动中，我们看见祂的能力常临在：当祂打雷、闪电、投掷霹雳，或当祂使一切重放光明时。如果你看不见天主，不应感到惊奇。风和风暴的冲击推动、摇撼、搅动万物，然而风和风暴却不进入我们的视线。同样，我们不能直视太阳，它是所有受造物看见的原因：瞳孔避开它的光线，观看者的目光模糊；如果看得太久，所有视觉能力都会熄灭。什么！你难道能承受太阳本身的建造者，光明的源泉本身，当你避开祂的闪电，躲开祂的雷霆时？你难道想用肉体的眼睛看见天主，而你自己连使你有生命、能说话的灵魂本身既不能看见也不能把握吗？而且，据说天主对人所作的事一无所知；祂居在天上，既不能察看一切，也不能认识个体。人啊，你错了，被欺骗了；因为天主何时远离我们呢？当天上、地上以及宇宙这个领域之外的一切都为天主所知、充满天主时。祂无处不在，不仅非常接近我们，而且注入我们内。因此，再看太阳：它固定在天上，却平等地散布在一切土地上；它临在于各处，与万物结合并混合；它的光辉从未受损坏。更何况创造了万物、察看万物的天主，对于祂没有什么是隐秘的，祂临在于黑暗中，临在于我们的思想中，如同在深沉的黑暗中。我们不仅在祂内行动，而且，我几乎可以说，我们与祂一同生活。\par

% structure: p0067 source=h2 target=section reason=relative-heading-rank; base=h2

\section{第三十三章 论点：即使有人说天主对犹太人毫无帮助，但\WorkTitle{犹太编年史}的写作者是最充分的见证，证明他们在被天主抛弃之前已经先离弃了天主。}

我们也不要因我们人数众多而自夸。我们自以为是许多，但在天主看来，我们非常少。我们区分民族和国家；对天主而言，整个世界是一个家庭。君王只能通过臣仆的侍奉知晓他王国的一切事务；天主教不需要情报。我们不仅在祂眼中生活，而且在祂怀中生活。但有人\Emph{反驳}说，犹太人自己以最大的迷信敬拜独一天主，有祭台和圣殿，这对他们毫无益处。如果你回忆后期的事件，却忘记或不知道早期的事件，你就犯了无知的罪。因为他们自己，当他们以贞洁、纯洁和虔诚敬拜我们的天主（祂是所有人的同一天主）时，当他们服从祂有益的诫命时，从少数变成无数，从贫穷变为富有，从奴隶变为君王；少数人压倒多数人；手无寸铁的人压倒武装的人，后者从他们面前逃跑，他们因天主的命令追赶，并有自然元素为他们效力。仔细阅读他们的圣经，或者如果你更喜欢罗马的著作，可以在\Person{弗拉维乌斯·约瑟夫斯}或\Person{安东尼·朱利安努斯}的书中（不提古代文献）查询关于犹太人的事，你就会知道，由于他们的邪恶，他们应得这样的命运，而且所发生的事没有一件不是事先预言给他们的，如果他们坚持他们的顽梗。因此你会明白，他们在被抛弃之前已经先离弃了，而且他们并不是像你亵渎所说的那样与他们的天主一同被俘，而是被天主作为祂纪律的逃兵所抛弃的。\par

% structure: p0069 source=h2 target=section reason=relative-heading-rank; base=h2

\section{第三十四章 论点：此外，这个世界将被火焚烧，这丝毫不值得惊奇，因为一切有开始的事物也有终结。古代哲学家并不反对世界可能被焚烧的观点。然而，显然天主从虚无中造人，也能将人从死亡中复活到永生。整个自然界也暗示着未来的复活。}

再者，关于世界被焚毁，认为火不会以意外的方式降在世界之上，或世界不会因此毁灭，是一种庸俗的错误。因为哪一个智者会怀疑，哪一个无知的人不知道，凡有开端的事物都会消亡，凡被造的都终有结束？天及其包含的一切，也必如它开始一样止息。海洋从甘泉中得到的滋养，将化为火的能力。斯多葛派一直相信，当湿气枯竭时，整个世界都将起火；伊壁鸠鲁派对元素燃尽和世界毁灭也持完全相同的观点。\Person{柏拉图}说，世界的各部分时而遭水淹，时而在交替变化中被烧毁；尽管他说世界本身被造得永恒而不灭，但他又补充说，对于天主自己——唯一的工匠——它既可分解又终有一死。因此，那团物质若被其建造者所毁，也不足为奇。你们看到，哲学家们所争论的正是我们所说的同样事情，并非我们在追随他们的足迹，而是他们从先知的神圣宣告中，模仿了那被败坏的真理的影儿。同样，最杰出的智者——首先是\Person{毕达哥拉斯}，主要是\Person{柏拉图}——以败坏而分裂的信仰传授了复活的教义；因为他们认为，身体分解后，唯有灵魂永远存留，并且常常转向其他新的身体。他们还以此为歪曲真理的方式补充说，人的灵魂会返回牛、鸟和野兽之中。毫无疑问，这种意见不值得哲学家探讨，只配小丑的粗俗笑谈。但就我们论证而言，这一点就够了：即使在这件事上，你们的智者也与我们有所和谐。然而，谁如此愚蠢或粗野，竟敢否认：人既能为天主首先形成，也能再被再造；死后归于无有，存在之前本是无有；既然从无中可以出生，那么从无中或许也可以恢复？而且，开始一个不存在的事物，比重复一个已存在的事物更为困难。你们认为，如果有任何东西从我们软弱的眼前消失，对天主而言就毁灭了吗？每一个身体，无论是化为尘土、融为湿气、压成灰烬，还是散成烟雾，都从我们面前消失，但被天主保存在元素的看管中。我们并不像你们所认为的那样，担心因丧葬而失去什么，而是采纳古老而更好的土葬习俗。因此，请看，为了我们的安慰，整个自然界都在预示着未来的复活。太阳落下又升起，星辰逝去又回来，花朵凋谢又重生，灌木在冬日凋零后重新长出叶子，种子若不腐烂就不能再发荣滋长：因此，坟墓中的身体恰如树木，在冬天用欺骗性的枯干隐藏其青翠。为何在冬天尚寒冷时，就急于让它复活并返回？我们也必须等待身体的春天。我也知道，许多人在自知其所应得时，与其说相信，不如说渴望死后归于虚无；因为他们宁愿完全熄灭，也不愿为了惩罚而被恢复。他们的错误也因今生所赐的自由和天主极大的忍耐而加深，天主的审判越是迟延，就越是公正。\par

\EditorialNote{\Quote{……从无} 从本章（若非其他章节）中，有人曾轻率地断言我们的作者认为灵魂与身体一同消亡，并要从无中被更新。据我理解，这一论证完全是针对对方而发，并未从作者自身视角陈述任何观点。他只给出了对其论证足够的东西，并不声称更多。他并非神职人员，他的作品也不是对信友的讲道。他挑战任何人否认：如果天主能从无中造人，他也能从无中重新造人。论证的其余部分是对物质不朽性的杰出断言，其措辞足以令现代科学满意；其含义是，灵魂在天主眼前并不比汽化并保存在元素看管中的身体更易消亡。}\par

% structure: p0072 source=h2 target=section reason=relative-heading-rank; base=h2

\section{第三十五章。论证：义人及虔诚者将获永福，不义者将受永罚。基督徒的道德远比外邦人圣洁。}

然而，人们在最有学识的诗人的书籍和诗歌中受到告诫，提到那火河，以及从\Place{斯堤克斯沼泽}以多种曲折流出的热流——这些事物是为永罚预备的，他们通过魔鬼的信息和其先知的训示得知，并传给了我们。因此，在他们当中，甚至连\Person{朱庇特}王也虔诚地指着灼热的河岸和漆黑的深渊起誓；因为他预知等候他和他的崇拜者的惩罚，而战栗。这些痛苦既无量度也无终止。那里，有灵性的火烧灼肢体又使之复原，以肢体为食又滋养它们。正如雷霆之火击中身体却不焚毁；正如\Place{埃特纳火山}和\Place{维苏威火山}之火，以及各处燃烧之地的火焰，燃烧却不耗尽；同样，那惩罚之火不是靠燃烧者的消耗来喂养，而是靠对其身体永无休止的侵蚀来滋养。但不认识天主的人，作为不敬神和不义者理应受折磨，除了不虔敬的人，无人怀疑去相信，因为不认识万物之父母和万物之主，并不比冒犯他们罪恶更小。尽管不认识天主已足以受罚，正如认识他有益于宽恕，然而如果我们基督徒与你们相比，尽管在某些方面我们的纪律较差，但我们会显得比你们好得多。因为你们禁止通奸，却仍犯通奸；我们生为\Emph{男子}只为了我们自己的妻子：你们惩罚已犯的罪行；对我们来说，甚至想到犯罪就是犯罪：你们害怕那些知道你们所作所为的人；我们甚至唯独害怕自己的良心，没有良心我们不能存在：最后，你们的数量使监狱爆满；但在那里没有基督徒，除非他因自己的宗教被控告，或是叛教者。\par

% structure: p0074 source=h2 target=section reason=relative-heading-rank; base=h2

\section{第三十六章。论证：命运无非是天主之安排。人的心智是自由的，因此其行动也是自由的：他的出身不被带入审判。基督徒因贫困而受责难，这不是耻辱而是荣耀；他们遭受身体的痛苦，这不是惩罚而是操练。}

谁也不可从源自\OriginalTerm{命运}{fatum}的事得安慰，也不可为它辩解。任凭所发生之事系于\OriginalTerm{幸运}{fortuna}的\OriginalTerm{定数}{dispositio}，然而心灵是自由的；因此被审判的是人的行为，而非人的尊严。因为命运若非天主对每个人所说的话，又是什么？祂既能预见我们的体质，便按各人的功过与品性，也为我们定下命运。因此在我们这里，受惩罚的不是我们出生时的星辰，而是我们性情的特殊性质受责备。关于命运，说得已经够多；或者，若因时间而说得太少，他日我们将更充分、更详尽地讨论此事。但我们许多人被称为贫穷，这并非我们的耻辱，而是我们的光荣；因为正如我们的心灵被奢侈所松弛，也被节俭所加强。然而，一个人若无所欠缺，不贪求他人之物，对天主富有，谁又能说他是贫穷呢？那才是真正的贫穷，即虽拥有许多，却渴望更多。但我仍将按照自己的感受来说。没有人生来就贫穷如斯。飞鸟无需祖产而活，牲畜日日得喂养；然而这些受造物乃为我们而生——所有这些，我们若不贪恋，便已拥有。因此，如行路者，行走愈轻省，愈为幸福；同样，在这人生旅途中，高举自己行于贫穷，不为财富的重担所气喘者，更为幸福。然而，纵使我们以为财富有益，也当向天主祈求。祂拥有万有，自能稍加满足我们；但我们宁愿轻看财富，也不愿拥有它们：我们更渴望纯洁，更祈求忍耐，宁可为善，也不愿挥霍；并且，我们感受到并忍受身体的人性灾祸，不是惩罚，而是争战。因为刚毅因软弱而加强，灾难常常是德行的操练；此外，心灵和身体的力量若无劳动的操练，便会萎靡。因此，你们所引以为榜样的所有强者，都因他们的磨难而辉煌兴盛。同样，天主既不能无力援助我们，也不会轻视我们，因为祂既是万民的统治者，又是爱护自己子民者。但在逆境中，祂察看并探询每一个人；祂在危险中衡量每个人的性情，直至最后的死亡；祂审察人的意志，深知在祂面前没有东西会灭亡。因此，如同金子经火，我们也由关键时刻显明出来。\par

% structure: p0076 source=h2 target=section reason=relative-heading-rank; base=h2

\section{第三十七章 论点：为宣认基督之名而最不公正施加的酷刑，是堪当天主的景象。在异教徒中一些最勇敢者与圣殉道者之间进行的比较。他宣称基督徒不参加公共表演和游行，因为他们极其确知这些活动既不虔敬又残忍。}

对于天主来说，当基督徒与痛苦作战时，这景象是何等美丽！当他被引向面对威胁、刑罚和酷刑；当他嘲笑死亡的喧嚣，践踏刽子手的恐怖；当他在君王和元首面前高举起自己的自由，只屈服于他所归属的天主；当他胜利而凯旋，践踏那对他宣判的人！因为他已经得胜，取得了所争夺的。哪位士兵在自己将军眼前不会更大胆地冒险呢？因为没有人能在试炼前得到赏报，然而将军不能给予他没有的：他不能保全生命，但他能使战争光荣。但天主的士兵在苦难中既不被遗弃，也不会因死亡而终结。因此，基督徒看似可怜，但实际上不可能真如此。你们自己把不幸的人捧上天；例如\Person{穆奇乌斯·斯凯沃拉}，当他意图行刺国王失败后，除非牺牲右手，否则就会死在敌人手中。而我们中有多少人忍受的不仅是右手，而是全身被焚烧，毫无痛苦的哭喊，尤其当他们本可以获释时！我能拿人们与\Person{穆奇乌斯}、\Person{阿基里乌斯}或\Person{勒古鲁斯}相比吗？然而我们中男孩和年轻妇女，以受启发的忍耐力，轻蔑地对待十字架、酷刑、野兽和一切可怕的刑罚。而你们这些可怜的人，难道不明白，没有人会无缘无故地甘愿受罚，也没有人能在没有天主的情况下忍受酷刑吗？除非，或许事实欺骗了你们，那些不认识天主的人富有、荣耀洋溢、权力超群。\Emph{可怜的人！}在这方面，他们被抬高得越高，是为了跌得更低。因为他们被养肥，作为惩罚的牺牲品，作为祭品被加冕以待宰杀。因此在这方面，有些人被提升到帝国和统治地位，以便权力的放纵行使能使他们的精神在毁灭灵魂特有的放纵中成为交易。因为，除认识天主之外，既然死亡必至，哪里会有坚实的幸福呢？幸福像梦一样，在抓住之前就溜走了。你是一位君王吗？然而你害怕的程度与别人怕你的程度相等；无论你身边有多少随从，在危险面前你仍是孤独的。你富有吗？但命运不可靠；携带大量行李，短暂的人生之旅不是被装备，而是被拖累。你夸耀\OriginalTerm{权杖}{fasces}和官袍吗？这是人虚妄的错误，是空洞的尊严崇拜——外表闪耀\OriginalTerm{紫袍}{purple}，内心却污秽。你因出身高贵而自高吗？你赞美你的父母吗？然而我们所有人命中注定的出生相同，只有德行才使我们不同。因此，我们这些以品行和谦逊为人所知的人，合理地戒绝邪恶的享乐，以及你们的炫耀和表演，我们知晓其与神圣事物相关的起源，并谴责其有害的诱惑。因为在战车比赛中，谁不因人群相互争吵的疯狂而战栗？或是在角斗游戏中教导谋杀？在戏剧表演中，疯狂同样不小，但放荡更加持久：因为此刻一个模仿者要么解说要么展示奸淫；此刻一个软弱的演员，一边假装情欲，一边暗示它；同一个演员羞辱你们的神明，将奸淫、叹息、仇恨归于它们；同样的演员用虚假的痛苦和空洞的姿势表情引动你们的眼泪。因此，你们实际上要求谋杀，同时在虚构中为它哭泣。\par

% structure: p0078 source=h2 target=section reason=relative-heading-rank; base=h2

\section{第38章 论点：基督徒避免与偶像祭品有关的事物，以免有人认为他们要么向魔鬼屈服，要么以宗教为耻。他们确实不鄙视花的颜色和香气，因为他们习惯于松散而随意地撒花，也用花环缠绕颈项；但给尸体戴上冠冕，他们认为多余而无用。此外，他们以同样的平静生活并埋葬死者，怀着非常确定的希望等待永恒幸福的华冠。因此，他们的宗教摒弃外邦人的一切迷信，应当被所有人接纳为真宗教。}

但我们轻视祭献的残余和奠酒所用的杯子，这并非承认恐惧，而是宣告我们真正的自由。因为尽管没有任何由天主作为不可侵犯的恩赐所造之物会被任何力量败坏，但我们仍然戒绝，以免有人以为我们是在向接受奠酒的魔鬼屈服，或是羞于我们的宗教。然而，有谁会怀疑我们在春天的花朵中放纵自己？我们同时采摘春天的玫瑰和百合，以及花朵中其他颜色怡人、气味芬芳之物。我们既松散自由地使用它们，也用它们编成花环缠绕颈项。请原谅我们，诚然，我们不戴花冠；我们习惯于用鼻孔吸入甜花的芳香，而不是用后脑勺或头发来吸气。我们也不为死者加冠。在这方面，我更加惊讶于你们，你们竟然将火炬献给一个没有生命或没有感觉的人，或将花环献给一个闻不到它的人——因为若是幸福的人，他不需要花；若是痛苦的人，他也不会从花中得到快乐。尽管如此，我们仍以生活中同样的宁静来装饰我们的葬礼；我们不系一朵枯萎的花环，而是佩戴一朵来自天主的、以永恒花朵编织的活花环，因为我们在天主的慷慨中既节制又安稳，因祂现世尊荣的确信而备受鼓舞，盼望未来的幸福。如此，我们得以在幸福中复活，并已经活在默想未来之中。那么就由\Person{苏格拉底}那个\Place{雅典}的小丑去理会吧，他承认自己一无所知，却夸耀于一个极其诡诈的恶魔的见证；也让\Person{阿尔克西劳}、\Person{卡涅阿德斯}、\Person{皮浪}，以及所有的\OriginalTerm{学院派哲学家}{Academic philosophers}都来斟酌吧；也让\Person{西摩尼得斯}永远推迟他意见的决定。我们鄙视哲学家们紧锁的眉头，我们知道他们是败德者、奸淫者、暴君，并且总是雄辩地谴责自己的恶习。我们这些不将智慧穿在衣饰上、而是藏在心中的人，我们不谈论大事，而是活出大事；我们夸耀我们已获得了他们以最大的热忱寻求却未能找到的东西。我们为何忘恩负义？我们为何嫉妒，如果神性的真理在我们这时代已经成熟？让我们享受我们的福祉，让我们以正直节制我们的判断；让迷信被约束；让不敬虔被赎除；让真宗教得以保存。\par

% structure: p0080 source=h2 target=omitted reason=duplicate-page-title

当\Person{奥克塔维乌斯}结束他的演讲时，我们一时间陷入沉默，凝神屏息；而我则沉浸在极大的惊叹之中，因为他竟能用论据、例证和阅读中得来的权威，如此完美地修饰那些本更容易感受而非言说的事物；并且他以哲学家们自己武装的武器击退了那些恶意的反对者，还揭示了真理既容易领会又令人愉悦。\par

% structure: p0082 source=h2 target=omitted reason=duplicate-page-title

因此，当我正默然心中思量这些事时，\Person{凯基利乌斯}突然说道：\Quote{我同样衷心地祝贺我的\Person{奥克塔维乌斯}和我自己，为我们所生活的这份安宁，并且我不等待决断。就这样我们得胜了：我并非不公正地将胜利归于自己。因为正如他是我的征服者，我也战胜了错误。因此，就问题的实质而言，我既承认天主眷顾，也向天主屈服；我同意我所拥有的生活方式的真诚。但仍有某些事留在我心中，并非抵触真理，而是为完满的训练所必需；明天，既然太阳已西斜，我们将以更合适和更从容的方式详细探讨。}\par

% structure: p0084 source=h2 target=omitted reason=duplicate-page-title

\Quote{但就我而言，}我说，\Quote{我更完全地为我们所有人而欢喜；因为\Person{奥克塔维乌斯}也为我赢得了胜利，既然评判的极端不公已从我身上移去。我也无法用赞美来表达他话语的价值：人，且仅是一个人的见证是软弱的。他从天主那里获得了光辉的赏报，蒙祂启示而辩护，靠祂助佑而取胜。}\par

此后，我们欣然离去：\Person{凯基利乌斯}因他已相信而欢喜；\Person{奥克塔维乌斯}因他已成功而欢喜；而我则因前者已相信、后者已获胜而欢喜。\par

\SourceRecord{Translated by Robert Ernest Wallis.；Ante-Nicene Fathers；Vol. 4.；Edited by Alexander Roberts, James Donaldson, and A. Cleveland Coxe.；Buffalo, NY: Christian Literature Publishing Co.,；1885.；<http://www.newadvent.org/fathers/0410.htm>.}
